\chapter{Claude Monetize's impressions of Noise}

\vspace{-\baselineskip}
\lettrine{C}{laude} Monetize was a man burdened by the weight of legacy. His granduncle, the legendary Claude Shannon, had given the world the mathematical foundations of \textbf{Information theory}, a framework that transformed how we think about data, noise, and communication. 

\begin{figure}[H]
\includegraphics[width=1.0\linewidth]{Illustrations/vign-claude-commonRoom.png}
\caption{Claude deaf to his Uncle's noise}
\end{figure}

Claude, instead had spent his early years navigating the murky waters of financial derivatives, analyzing reinsurance risks for a large corporation that seemed to delight in avoiding payouts. Lucrative work indeed, which left him bereft of meaning and a yearning for a legacy of Uncle Claude proportions. So of late, Claude had been drawn back to pure mathematics, searching for solace within the entropy of prime numbers through uncle Shannon’s signal-to-noise ratio formula. Eleanor’s recent call (that he had received in the post), for a circle offered him a chance to work with others who shared his lofty aspirations.

\smallskip

Claude often found himself sketching out Shannon’s formula, which for \(C\) the channel capacity (bits per second), \(B\) the bandwidth (hertz), \(S\) the signal power (watts), and \(N\) the noise power quantifies the capacity of a communication channel as a function of the signal-to-noise ratio:
\begin{align}
    C = B \log_2 \left(1 + \frac{S}{N}\right).
\end{align} \label{eq:CludeSignalNoise}
For Claude, this offered a way to think about the humble number line, where the signal (\textbf{semi-prime} numbers, meaningful patterns) were hidden in the noise of \textbf{composite} numbers. As such he had in mind the eponymous \textit{Shannon entropy}, \(H(X)\) the measure of uncertainty in a probability distribution of the random variable \(X\):
\[
H(X) = - \sum_i P(x_i) \log_2 P(x_i),
\]
where\(P(x_i)\) is the probability of outcome \(x_i\) and Einstein conventions means we sum, \( \sum_i\) over all such i. This Entropy relationship reminded Claude of the prime number theorem, which describes the distribution of primes using the approximation:
\[
\pi(x) \sim \frac{x}{\log x},
\]
where: \(\pi(x)\) is the number of primes less than or equal to \(x\) and  \(\log x\) is the natural logarithm of \(x\).



Claude saw a conceptual overlap between entropy and the attenuation of primes along the number line. Just as entropy measures the uncertainty of a message, \(\pi\sim\frac{x}{\log x}\) measures how sparse primes become as \(x\) grows. Claude wondered: could the tools of information theory extract meaningful patterns from the "prime noise"?

\begin{figure}[H]
\centering
\includegraphics[width=0.6\linewidth]{Illustrations/vign-claudeHotel.png}
\caption{Claude not the worst for wear for his financial indulgences}
\end{figure}
\smallskip

\textbf{Semiprimes as Moiré Interference with Angular Modulation:}

\noindent
To this end, Claude was struck by the caustic complexity in Moiré’s explorations of interference patterns, especially those emerging from angular modulations that give rise to spiraling, dynamic meshes. Accordingly he defined two Fourier \(\pi(r)\)-modulated fields in polar coordinates as:
\begin{align}
f_1(r,\theta) = \sin\left( \frac{r}{\log r} \right), \quad
f_2(r,\theta) = \sin\left( \frac{r}{\log r} + \omega(\theta + \delta) \right),
\end{align} \label{eq:Moire}

where $\omega$ is an angular frequency constant and $\delta$ is a small angular offset which are used to introduce an angular interference, producing the rotating spiral-like Moiré mesh.

\begin{figure}[H]
\centering
\includegraphics[width=0.7\linewidth]{Illustrations/vign-shannonMoire (1).png}
\caption{Claude imagining Moire patterns}
\end{figure}

The moiré pattern emerges from the product of the two fields:
\[
M(r, \theta) = f_1(r, \theta) \cdot f_2(r, \theta + \delta).
\]
 \begin{figure}[H]
\centering
\includegraphics[width=1.0\linewidth]{Illustrations/moire_collage_with_theta.png}
\caption{\href{https://colab.research.google.com/drive/1MVE7qkP1DQFbqvZPLsoTsqvqHJBitpa7?usp=sharing}{Moire Mesh} with highlighted semi-primes}
\end{figure}
Accordingly, Semiprimes (numbers of the form $pq$ for distinct primes $p \ne q$), Claude conjectured would align with constructive interference zones of $M(r, \theta)$, appearing as ridges in its contour map:
\[
\text{Semiprimes} \approx \left\{ (r,\theta) \,\middle|\, \nabla M(r,\theta) = 0 \text{ and } M \text{ locally maximal} \right\}.
\]

\newpage
\section*{A Reptend Debate that will go on and on}

Eleanor sweeps into Claude's office unannounced, as was her habit, carrying a rolled-up poster of what appeared to be interference patterns from her lab's latest acoustic waveguide experiments.

``Still playing with your Moiré spirals and pretending they reveal the secrets of semiprimes?'' she asked, not unkindly, but with that edge of skepticism Claude had come to expect.

He didn't look up from his screen, where the scatter plot from 

\noindent
\href{reptend_capacity_function.html}{\texttt{reptend\_capacity\_function.html}} showed the distribution of $\rho(p) = \varphi(p-1)/(p-1)$ across the first thousand primes. ``And you're still pretending that physical intuition trumps rigorous proof?''

``Touché Claude, touché.'' Eleanor unrolls her poster across his desk, obscuring his keyboard. ``But look: these acoustic standing waves in our resonator cavity form hexagonal patterns at certain frequencies. The nodes align precisely where I'd expect Eisenstein primes to cluster if we map complex plane to physical space. Nature \emph{knows} about the hexagonal lattice structure before we prove anything about it.''

Claude finally looks up, pushing his glasses further down his nose to accentuate his sniffiness and simmering contempt for all her enthusiasm. ``Dear Eleanor, nature doesn't \emph{know} anything. It simply \emph{is}. Your resonator follows Maxwell's equations, which are deterministic and conformally invariant te boot. Any correspondence to number-theoretic structures is an artefact of your modelling or worse, you're unconsciously selecting experimental parameters to match your expectations.''

``Am I?'' She tapped the poster. ``Then explain why the interference maxima occur at radii that scale as $r \sim p^{1/2}$ where $p$ are the reptend primes. That's not in my experimental design ... that's emerged from the physics.''

Claude arching way from her, but intrigued despite himself. ``Show me the data.''

Eleanor pulled out her tablet, displaying a log-scale plot. ``Resonance peaks. I'm measuring acoustic pressure amplitude at different radial distances from the center of a circular membrane. Here ...'' she zoomed in, `` ... these peaks correspond to $r = \sqrt{7}, \sqrt{17}, \sqrt{19}$\ldots''

``The base-10 reptend primes,'' Claude murmured... He opened his recent paper on ResearchGate: \emph{From Reptend Primes to Fermat Primes: A Base-Invariant Characterization via the Reptend Capacity Function} \cite{mcculloch2025reptend}. ``But Eleanor, reptend status is base-dependent. Your membrane doesn't care about base-10. It's not counting in decimal.''

``Exactly!'' Her eyes lit up. ``That's what I'm trying to tell you. My membrane is responding to something \emph{base-invariant} about these primes. And your paper ...'' she grabbed his mouse and scrolled to Section 3, ..., ``here, it introduces exactly the right concept. The reptend \emph{capacity} function, $\rho(p) = \varphi(p-1)/(p-1)$. That's not base-dependent.''

Claude felt the uncomfortable sensation of having his own work used against him. ``The capacity function measures the \emph{fraction} of bases in which $p$ is reptend. It's still fundamentally about representation systems, about information encoding. What could a membrane possibly care about that?''

``Information \emph{transmission},'' Eleanor countered. She pulled up Shannon's formula, (\ref{eq:CludeSignalNoise}) on her own tablet. ``You love this equation: channel capacity as a function of signal-to-noise ratio. Now, what if the `noise' in the acoustic system isn't random? What if it's structured by the composite numbers, and the `signal', that is the coherent standing waves, preferentially amplifies at positions corresponding to primes with high reptend capacity?''

``That's\ldots'' Claude paused. ``That's not completely insane, actually. But you'd need to show me the mechanism. Why would $\rho(p)$ couple to acoustic resonance?''

Eleanor sat down, suddenly serious. ``I don't know yet. But consider: primes with high reptend capacity have $p-1$ with simple factorization, ... ideally just $p-1 = 2^k$ for Fermat primes, where $\rho(p) = 1/2$ exactly. That's maximal \emph{periodicity} in the multiplicative group structure. And what are standing waves if not periodic structures?''

``You're conflating algebraic periodicity with physical periodicity.''

``Am I? Or are they different manifestations of the same underlying principle?'' She pointed to his Moiré interference equations, \ref{eq:Moire}:

``You're using $r/\log r$ explicitly because of the prime number theorem, $\pi(x) \sim x/\log x$. You're already assuming physical space can encode number-theoretic information. How is my acoustic cavity any different?''

Claude stood, pacing to the window. ``The difference, Eleanor, is that my Moiré patterns are \emph{models},... mere mathematical analogies. I'm not claiming the prime numbers exist as physical interference patterns in spacetime. I'm using the formalism of wave interference as a \emph{metaphor} for understanding the multiplicative structure of integers.''

``Are you sure about that?'' Eleanor joined him at the window, looking out over the quad where two students were chalking partial differential equation puns onto the pavement in a misguided attempt at public mathematics. ``Because I've been reading your paper carefully. This business with Fermat primes achieving maximal reptend capacity. Here this: $\rho(p) = 1/2$ exactly when $p-1 = 2^k$, that's not just numerology. There's real structure there.''

``Of course there is. It's a theorem.'' Claude pulled it up :

\begin{theorem}[Characterization of Maximal Capacity]
For odd prime $p$, the following are equivalent:
\begin{enumerate}
\item $\rho(p) = 1/2$ ,\quad $\varphi(p-1) = (p-1)/2$
\item $p-1 = 2^k$ for some $k \geq 1$ ,\quad
 $p$ is a Fermat prime
\end{enumerate}
\end{theorem}

``This is rigorous mathematics,'' he said. ``We proved this using Fermat's little theorem. It has nothing to do with membranes or waveguides or your hexagonal obsessions.''

Eleanor smiled, that infuriating, knowing smile. ``And yet your proof relies on the fact that $\varphi(2^k) = 2^{k-1}$, which is a statement about \emph{binary structure}. Binary, Claude. Powers of two. Octave relationships. Harmonic resonance.''

``No I am not having that. That really is just numerology, not mathematics.''

``Is it?'' She pulled up another plot. ``This is the surprise metric from your paper: $S(p) = -\log_2|\rho(p) - 1/2|$ for near-maximal primes. You're measuring how `unexpectedly close' a prime comes to Fermat-like behavior in \emph{bits}. Bits! Shannon entropy hiding in plain sight.''

Claude opened his mouth to object, then paused. She had a point. He'd been thinking about the surprise metric as purely mathematical. Rather as a quantification of deviation from the maximum. But Eleanor was right: it had the structure of information content, of \textit{surprisal} in the Shannon sense.

``Go on,'' he said carefully.

Eleanor pulled out a chair, sensing an opening. ``Your paper proves that Fermat primes ... and only Fermat primes ... achieve maximum reptend capacity of $\rho(p) = 1/2$. That's a \emph{physical} maximum, Claude. It's not some arbitrary mathematical boundary emerging because $\varphi(2^k) = 2^{k-1}$, and that factor of one-half appears everywhere in physics.''

``Give me an example,'' Claude challenged.

``\textit{Nyquist-Shannon} sampling theorem,'' Eleanor shot back immediately. ``You need to sample at twice the maximum frequency to perfectly reconstruct a signal. The factor of 2. Your Fermat primes achieve $\rho(p) = 1/2$ because $p-1 = 2^k$, ... they're \emph{balanced} between structure and randomness in the same way that critical sampling in signal processing sits exactly at the Nyquist frequency. It's the same $1/2$ showing up. That can't be coincidence.''

Claude oohing his mouth to object, closed it, instead pulling up the \href{exact_vs_approximate_50_percent.html}{\texttt{exact\_vs\_approximate\_50\_percent.html}} tool, showing primes that approached but didn't achieve the 50\% maximum. ``The primes with $p-1 = 2q$ where $q$ is large\ldots\ they approach 50\% asymptotically. Like $(q-1)/(2q) \to 1/2$.''

``Exactly like a resonance approaching a critical point,'' Eleanor said quietly. ``Claude, what if the reptend capacity function is telling us about actual \emph{information capacity} in the distribution of primes? Your Shannon entropy formula, ... '' she pointed to his whiteboard where $H(X) = -\sum_i P(x_i) \log_2 P(x_i)$ was still written, ``, measures uncertainty. And the prime number theorem says primes become sparser as $\pi(x) \sim x/\log x$. That logarithm is doing \emph{work}. It's not decorative.''

Claude stood and walked to the window and idly looked out thoughtlessly. ``You're saying there's a physical interpretation of why Fermat primes achieve maximum reptend capacity? That $\rho(p) = 1/2$ for $p = 2^{2^n} + 1$ isn't just an arithmetic curiosity?''

``I'm saying your Moiré patterns, ...'' Eleanor gestured at his abandoned figures on the desk, ``..., might be more right than you think. When $p-1 = 2^k$, you have pure power-of-two structure. No odd prime factors to create `interference.' The signal is cleanest. The information capacity is maximal.''

``But that's\ldots\ that's not a proof of anything.''

``No,'' Eleanor agreed. ``But it's a \emph{direction}. Your paper proves Fermat primes are unique in achieving $\rho(p) = 1/2$. That's beautiful mathematics. But \emph{why} should primitive root density matter? Why should the structure of $(\mathbb{Z}/p\mathbb{Z})^*$ connect to anything physical?''

Claude turned back to her. ``Because cyclotomic fields appear in physics. Crystallographic restriction theorem. Quasicrystals. The Eisenstein integers you mentioned: they describe the hexagonal lattice because $\omega^3 = 1$, and that rotational symmetry is the same symmetry you see in your acoustic resonator.''

``Now you're getting it.'' Eleanor smiled. ``The mathematics isn't \emph{just} reflecting abstract structure. It's reflecting physical constraints on how information can be organized in space. Your base-invariant characterization, ...'' she tapped his paper again, ``... is picking out primes whose multiplicative structure is maximally `simple' in a precise information-theoretic sense. And my experiments suggest that physical systems naturally select for that simplicity.''

Claude sat back down, reopening \href{base_invariant_reptends.html}{\texttt{base\_invariant\_reptends.html}}, his interactive tool showing which primes were reptend across multiple bases. ``The Fermat primes appear as bright spots across all bases\ldots''

``Because they're fundamental,'' Eleanor said. ``Not just mathematically. \emph{Fundamentally}. Claude, you spent years analyzing financial noise, trying to extract some signal. Maybe the universe has the same problem. And maybe, ... '' she pointed at $\rho(p) = 1/2$, `` ... this is telling us something about how the universe solves it.''

Claude was silent for a long moment. Then: ``Eleanor, this is wildly speculative.''

``Of course it is. That's what makes it interesting.''

Claude pulled up the surprise metric tool, showing primes that \emph{almost} achieved 50\% capacity. ``These primes with high surprise scores\ldots\ if there's a physical interpretation, what are they?''

``Near-critical states?'' Eleanor suggested. ``Systems that are \emph{almost} at a phase transition but not quite? Claude, I don't know. But I know that when mathematics and physics both point at the same structure from different directions\ldots'' She trailed off.

``We should pay attention,'' Claude finished. He looked at his Shannon entropy formula, then at Eleanor's acoustic data, then back at the reptend capacity function. ``Even if it makes us uncomfortable.''

``Especially then,'' Eleanor said. ``The uncomfortable questions are usually the interesting ones.'' As she gathered her poster to leave, Eleanor paused. ``Claude? One more thing.'' ...``Yes?''

``Your paper notes that the reptend capacity function might be novel. But the primitive root density; that's classical. Ireland and Rosen, standard textbook material.'' ... ``I'm aware.''

``So what you've really done'' Eleanor's eyes gleamed, ``is taken something abstract and asked a \emph{physical} question about it. What does it \emph{mean} for a prime to be reptend in many bases? You framed it as base-invariance, like coordinate-invariance in physics. That's not pure mathematics anymore, Claude. That's mathematical physics.''

She left before he could protest, leaving Claude staring at his equations with an expression somewhere between annoyance and wonder.

On his whiteboard, he added one more line by the Shannon formula:
\[
\rho(p) = \frac{\varphi(p-1)}{p-1} \quad \text{???} \quad C = B \log_2\left(1 + \frac{S}{N}\right)
\]

The question marks stayed there for weeks.

\begin{center}
\decofourleft \quad \decofourright
\end{center}

\vspace{1em}


\section*{Timetabling of Girlfriends}

Claude finds himself sitting in a pub with young Lucas, whom Eleanor has sent to do her bidding, hoping her signal is not lost by the effervescent but overly verbose young gun. Claude is busy recalling his prime youthful days spent in Totnes, Devon. Leaning back in his chair, watching Lucas scribbling notes like he was recording the secrets of life, he recounts:

"I am reminded of my younger days living here in Totnes, midway between Exeter and Plymouth. Back then, I had two girlfriends: Sarah up in Exeter and Lucy down in Plymouth." Claude mused, “You know, Lucas, back in the day, I let the trains decide my love life. I thought I was being clever, but it turns out I was a a misguided fatalist.”  

“Trains? You mean like flipping a coin?” Lucas queried.

\smallskip

“Sort of,” Claude replied. “I’d go to Totnes station at random times. If the train heading to Exeter arrived first, I’d go see Sarah. If the Plymouth train came first, I’d visit Lucy. It seemed fair. But something strange happened.”  

“What was that?” Lucas pressed, leaning forward.

\smallskip

Claude playfully sighing. “Well, after a couple of months, I’d visited Lucy twenty-eight times and Sarah only twice. Naturally, I thought fate was pushing me toward Lucy, but then I figured out the truth.”  

Lucas pen poised over his notebook: “So, what was it?” he prodded.

“It was the train schedules,” Claude came clean. “The Exeter trains followed primes of the form \(4n+3\): they came at 7, 19, 31, and 43 minutes past the hour. The Plymouth trains, on the other hand, followed primes of the form \(4n+1\): 5, 17, 29, and 41 minutes past the hour. And here’s the key: the downflow to Plymouth always happened just before the upflow to Exeter.”  

\smallskip

“What difference does that make?” Lucas asked, now plainly feigning puzzlement to keep the codger sweet.

\smallskip

“Let’s say I stroll up to the station at a random time,” I explained. “Suppose it’s 6 minutes past the hour. The Plymouth train (5 minutes past) has just left, so I’ve got nearly 12 minutes to wait for the next one. But the Exeter train (7 minutes past) arrives almost immediately. On the other hand, if I arrive at 16 minutes past, the Plymouth train (17 minutes past) is due first, leaving the Exeter train (19 minutes past) slightly behind.

\smallskip

“In short, the Plymouth trains had just enough of a head start to make them more likely. Even though the gaps were identical, the sequence gave Lucy an unfair advantage. Timing, mate. Timing.”  

\smallskip

Lucas: “So it wasn’t fate. It was just... bad timing?” 

\smallskip

“Exactly,” I retorted. “It had nothing to do with Lucy being ‘meant to be.’ It was just that the Plymouth train schedule gave her the edge. Subtle, but it skewed the odds massively in her favor.”  

“And Sarah never figured it out?” Lucas inquired.

“Neither of them did,” I confessed. “But by then, the damage was done. Lucy thought I was too distracted, and Sarah thought I was unreliable. Totnes kicked me out metaphorically, and I learned a hard lesson: never let a timetable—or prime numbers—run your love life.”  

\begin{figure}[H]
\centering
\includegraphics[width=0.8\linewidth]{Illustrations/vign-train.png}
\caption{Claude hoping for Sarah}
\end{figure}

“So what’s the takeaway?” Lucas asked thoughtfully.

“The smallest differences can have the biggest effects,” I said, leaning back. And sometimes, what feels like fate is just a quirk of probability—or arithmetic.”  

“And don’t juggle two girlfriends at once?” Lucas grinned.
\newpage
\section*{Dirichlet on Primes in Arithmetic Progressions}
Consider the primes of the forms $4n+1$ and $4n+3$ in the range $[1, 100]$:
\begin{itemize}
    \item $4n+1$: $5, 13, 17, 29, 37, 41, 53, 61, 73, 89$ (10 primes).
    \item $4n+3$: $3, 7, 11, 19, 23, 31, 43, 47, 59, 67, 71, 79, 83, 97$ (14 primes).
\end{itemize}

Dirichlet's theorem states that for any two coprime integers $a$ and $d$, the arithmetic progression
\[ a, a+d, a+2d, \dots \]
contains infinitely many primes. Furthermore, the distribution of primes in different residue classes modulo $d$ is asymptotically uniform. For $d=4$, this implies:
\begin{itemize}
    \item The residue class $1 \pmod{4}$ corresponds to $a=1$ and $d=4$ which generates primes of the form $4n+1$, such as $5, 13, 17, \dots$.
    \item The residue class $3 \pmod{4}$ corresponds to $a=3$ and $d=4$ which generates primes of the form $4n+3$, such as $3, 7, 11, \dots$.
\end{itemize}

While these residue classes each contain infinitely many primes and are equally distributed in the limit as $x \to \infty$, local fluctuations in density can occur, which impact the likelihood of encountering primes of a specific type. This is due to:
\begin{itemize}
    \item interaction of congruence conditions with sieve of Eratosthenes.
    \item Statistical biases that favor certain residue classes in specific ranges.
\end{itemize}

\noindent
In the python \href{https://colab.research.google.com/drive/186PSbWnCrMV-TtRFedcnDVqhpmeYEqyH?usp=sharing}{coded} Monte Carlo simulation, we might consider composite numbers $c = p_1^{e_1} p_2^{e_2} \cdots p_k^{e_k}$ and step by $m = \frac{c}{d}$, where $d$ is a randomly chosen divisor. The nearest prime to $c \pm m$ is more likely to be of type $4n+3$ due to their greater local density. When moving forward from a composite number $c$, the nearest prime is determined by the first residue class modulo 4 that contains a prime. The first $4n+1$ primes will often appear slightly earlier than the corresponding $4n+3$ prime in many intervals analogous to a train schedule where trains arrive at $4n+1$: 5, 17, 29, 41 and $4n+3$: 7, 19, 31, 43 minutes past the hour. If a person arrives at a random time, they are almost just as likely to miss a train at $7$ past as at $5$ past, but the earlier train is more likely to be the next train caught. The histogram of nearest primes demonstrates the slight skew toward $4n+1$ primes.
\begin{figure}[H]
\centering
\includegraphics[width=0.6\linewidth]{Illustrations/histogram4n+1.png}
\caption{Nearest Primes ($4n+1$ vs $4n+3$)}
\end{figure}

\begin{figure}[H]
\centering 
\includegraphics[width=0.6\linewidth]{Illustrations/scattertrain.png}
\caption{$\log(m)$ vs $\log(c)$ of Nearest Prime}
\end{figure}

Both $4n+1$ and $4n+3$ classes have gaps where no primes exist in certain intervals: in the range $[90, 100]$, there are no $4n+3$ primes, while in the range $[120, 130]$, there are no $4n+1$ primes. 
The missing primes create irregularities in the distribution, making it possible for $4n+1$ primes to "fill in" certain intervals where $4n+3$ primes are absent.

The plot of composite numbers versus step sizes confirms this. We see that a purely computational approach falls short of finding a theorem such as Dirichlet's. In fact it may have led us to make some erroneous assertions akin to Lucy being favored over Sarah due to the alignment of the stars.

\noindent
\begin{minipage}[t]{0.4\textwidth}
  \vspace{0pt} % Alignment trick
The plot to the right is of the distribution of semi-primes up to 1000, where each point represents a semi-prime composed of two prime factors \( p \) and \( q \). The \( x \)-axis and \( y \)-axis correspond to the smaller and larger prime factors of the semi-primes, respectively.
  
\end{minipage}%
\hfill % Adds horizontal space between the minipages if needed
\begin{minipage}[t]{0.55\textwidth}
  \vspace{0pt} % Alignment trick
  \includegraphics[width=\textwidth]{Illustrations/bubbleSemiPrimes.png}
  \captionof{figure}{Totients of Semi-Primes}
  \label{fig:semiprimes}
\end{minipage}

The size of each point is proportional to the totient of the semi-prime, while the color indicates the nature of its Mobius function value - red for semi-primes with square factors and green for others.

% >>>>>> ANNEX B (cut-sheet v2): quadratic forms, prime products and residue chains — block commented out below, pointer left in text <<<<<<
\textit{(Quadratic forms, prime products and residue chains --- the full working is set out in Annexe B.)}

% \section*{Quadratic Forms and Prime Products}
% 
% The Prime numbers classifications based on their congruence modulo 4 play a fundamental role in number theory. If two primes $p$ and $q$ are of the form $4n+1$, then their product is also of the form $4n+1$:
% \begin{equation*}
% (4a+1)(4b+1) = 16ab + 4a + 4b + 1 = 4(4ab + a + b) + 1.
% \end{equation*}
% Consider $5$ and $13$, both of the form $4n+1$,  and their product can be written as a sum of two squares::
% \begin{equation*}
% 5 \times 13 = 65= 8^2 + 1^2.
% \end{equation*}
% Furthermore, by Fermat’s theorem on sums of two squares, every prime $p \equiv 1 \pmod{4}$ can be written as:
% \begin{equation*}
% p = a^2 + b^2.
% \end{equation*}
% Since the product of two sums of squares remains a sum of squares:
% \begin{equation*}
% (a^2 + b^2)(c^2 + d^2) = (ac - bd)^2 + (ad + bc)^2,
% \end{equation*}
% we conclude that the product of two $4n+1$ primes can always be written as a sum of two squares. If two primes $p, q$ are both of the form $4n+3$, their product is also of the form $4n+1$:
% \begin{equation*}
% (4a+3)(4b+3) = 16ab + 12a + 12b + 9 = 4(4ab + 3a + 3b + 2) + 1.
% \end{equation*}
% However, unlike the previous case, such a product cannot be written as a sum of two squares. The proof follows from considering some modular arithmetic. If a number is expressible as a sum of two squares,
% \begin{equation*}
% a^2 + b^2 \equiv 0, 1, 2 \pmod{4},
% \end{equation*}
% whereas $4n+3$ numbers satisfy 
% \(
% p \equiv 3 \pmod{4},
% \)
%  which is impossible since a sum of squares never yields 3 modulo 4. Thus, a product of two $4n+3$ primes can never be expressed as a sum of two squares. For example, consider $7$ and $11$, both of the form $4n+3$:
% \begin{equation*}
% 7 \times 11 = 77,
% \end{equation*}
% and yet though the composite formed of their product is $4n+1$, it cannot be written as a sum of two squares.
% While we will call $4n+1$ composites \textbf{Fermatiles}
% we will call those built of $4n+3$  primes as \textbf{Gausscomps} emphasizing their distinct number-theoretic behavior. 
% 
% \newpage
% 
% \subsection*{Quadratic residues}
% 
% A \textbf{Quadratic Residue} (QR) modulo \( p \), depends on the nature of p, being any integer for which there exists some \( m \) satisfying:
% \[
% m^2 \equiv r \pmod{p}.
% \]
% \begin{itemize}
%     \item If \( p \equiv 1 \pmod{4} \), then \textbf{\(-1\) is a quadratic residue}, meaning there exists some \( m \) such that:
%     \[
%     m^2 \equiv -1 \pmod{p}.
%     \]
%     This property leads to \textbf{Fermat’s theorem on sums of two squares}, ensuring any \( 4n+1 \) prime can be expressed as:
%     \[
%     p = a^2 + b^2.
%     \]
%     \item If \( p \equiv 3 \pmod{4} \), then \textbf{\(-1\) is not a quadratic residue}, meaning no integer \( m \) satisfies:
%     \[
%     m^2 \equiv -1 \pmod{p}.
%     \]
%     Consequently, primes of the form \( 4n+3 \) \textbf{cannot} be written as a sum of two squares.
% \end{itemize}
% 
% The Legendre symbol \( \left(\frac{a}{p}\right) \) offers a glimpse into the fundamental nature of QRs modulo an odd prime \( p \):
% \[
% \left(\frac{a}{p}\right) = 
% \begin{cases} 
% 1 & \text{if } a \text{ is a quadratic residue modulo } p, \text{ and } a \not\equiv 0 \pmod{p}, \\
% -1 & \text{if } a \text{ is a non-quadratic residue modulo } p, \\
% 0 & \text{if } a \equiv 0 \pmod{p}.
% \end{cases}
% \]
% Consider, for example, the evaluation of \( \left(\frac{23}{3}\right) \). Since \( 23 \equiv 2 \pmod{3} \), we reduce to checking \( \left(\frac{2}{3}\right) \). Squaring elements modulo 3:
% \[
% 1^2 \equiv 1 \pmod{3}, \quad 2^2 \equiv 4 \equiv 1 \pmod{3},
% \]
% reveals that 2 is \textit{not} a QR, thus \( \left(\frac{2}{3}\right) = -1 \), and therefore, \( \left(\frac{23}{3}\right) = -1 \). 
% 
% Similarly, computing \( \left(\frac{3}{23}\right) \) by testing powers of 3 modulo 23, we find:
% \[
% 3^{11} \equiv 1 \pmod{23},
% \]
% confirming that 3 is a QR modulo 23. That is, \( \left(\frac{3}{23}\right) = 1 \). Already, we begin to see an asymmetry between \( \left(\frac{p}{q}\right) \) and \( \left(\frac{q}{p}\right) \), a phenomenon captured in the law of quadratic reciprocity. To capture the broader structure of Legendre symbols, consider the matrix representation for small odd primes:
% \[
% \begin{array}{c|cccccccc}
% p \backslash q & 3  & 5  & 7  & 11 & 13 & 17 & 19 & 23 \\
% \hline
% 3  &  0  & -1 & -1 &  1 &  1 & -1 & -1 &  1 \\
% 5  & -1  &  0 & -1 &  1 & -1 & -1 &  1 & -1 \\
% 7  &  1  & -1 &  0 & -1 & -1 & -1 &  1 & -1 \\
% 11 & -1  &  1 &  1 &  0 & -1 & -1 &  1 & -1 \\
% 13 &  1  & -1 & -1 & -1 &  0 &  1 & -1 &  1 \\
% 17 & -1  & -1 & -1 & -1 &  1 &  0 &  1 & -1 \\
% 19 &  1  &  1 & -1 & -1 & -1 &  1 &  0 & -1 \\
% 23 & -1  & -1 &  1 &  1 &  1 & -1 &  1 &  0 \\
% \end{array}
% \]
% 
% \subsection*{Quadratic Reciprocity} 
% The \textbf{quadratic reciprocity} theorem states:
% \[
% \left(\frac{p}{q}\right)\left(\frac{q}{p}\right) = (-1)^{\frac{(p-1)(q-1)}{4}}.
% \]
% From this, we predict the following:
% \begin{itemize}
%     \item If \textbf{either} \( p \) or \( q \) is of the form \( 4k+1 \) (Fermat-type primes), then:
%     \[
%     \left(\frac{p}{q}\right) = \left(\frac{q}{p}\right),
%     \]
%     which implies a symmetric pattern in the heatmap.
%     
%     \item If \textbf{both} \( p \) and \( q \) are of the form \( 4k+3 \) (Gauss-type primes), then:
%     \[
%     \left(\frac{p}{q}\right) = -\left(\frac{q}{p}\right),
%     \]
%     leading to an \textbf{asymmetric flip} along the diagonal.
% 
% \end{itemize}
% 
% A heat map reveals the \textcolor{cyan}{\textbf{Gauss primes}} of the form \( p \equiv 3 \pmod{4} \), and \textcolor{pink}{\textbf{Fermat primes}}, satisfying \( p \equiv 1 \pmod{4} \).
% 
% \begin{figure}[H]\centering
% \includegraphics[width=0.8\textwidth]{Illustrations/legendreResidues100framed.png}
% \caption{Legendre Symbol Matrix: Fermat primes (\textcolor{pink}{pink}) vs. Gauss primes (\textcolor{cyan}{cyan}).}
% \label{matrixlegFrame:fig}
% \end{figure}
% \noindent
% \textbf{Key Observations from the Heatmap}:
% \begin{enumerate}
%     \item \textbf{Diagonal Consistency:} The self-pairing of primes (\( p=p \)) always yields \( 1 \), reinforcing that a number is always a residue of itself.
%     
%     \item \textbf{Symmetry for Fermat Primes (\( 4k+1 \))}:  
%     The pink-framed entries correspond to pairs where either \( p \) or \( q \) is of the form \( 4k+1 \). As expected, these sections mirror perfectly across the diagonal.
% 
%     \item \textbf{Flipped Signs for Gauss Primes (\( 4k+3 \))}:  
%     The cyan-framed entries correspond to interactions between two \( 4k+3 \) primes.  As predicted by reciprocity, these pairs show sign flips across the diagonal, appearing as negative reflections.
% \end{enumerate}
% 
% \bigskip
% Eleanor with due coaxing from Lucas decides to meet Claude. They are mid conversation already, and Claude is warming up nicely:
% 
% "The deeper we probe, the more striking the connections become. While the Legendre symbol was historically a tool for solving quadratic congruences, it has since found new life in the study of elliptic curves, modular forms, and the lattice-based cryptography that is your big thing now Eleanor?"
% 
% Eleanor isn't in that space. She has the priest on her mind. His words, mostly.
% 
% \begin{itemize}[leftmargin=2cm,labelsep=1em]
%     \item[\textbf{Eleanor:}] 
%     Claude, I have been speaking with a priest. \\
%     We traced the world as a Möbius strip: \\
%     Myth $\to$ Matter $\to$ Life $\to$ Mind $\to$ Myth. \\
%     A liturgical procession, ever circling, ever new.
% 
%     \item[\textbf{Claude:}] (changing gear effortlessly)
%     Hah! That procession sounds like redundancy in a code. \\
%     Repetition  that is not waste; but rather is error-correction, through mutation. \\
%     Liturgy keeps the signal strong despite the noise that comes with the passing of time.
% 
%     \item[\textbf{Eleanor:}] 
%     You mean the chants and beads are bits? \\
%     That faith is a, kinfd of, ..., checksum?
% 
%     \item[\textbf{Claude:}] 
%     Exactly. Every myth retold is like re-sending the packet. \\
%     Drop a bead, drop a syllable, the cycle still conveys the whole.
% 
%     \item[\textbf{Eleanor:}] 
%     But what of the mystery making, Claude? \\
%     The priest would not call it redundancy, but sacrament.
% 
%     \item[\textbf{Claude:}] 
%     Ah, but even sacrament has an information rate. \\
%     Too much compression, and you lose the soul of the signal.
% 
%     \item[\textbf{Eleanor:}] 
%     Then we agree: probability is not only epistemic, not only ontic, but a rhythm, ...,  a probabilitising procession. A rhythm that keeps meaning alive:
%     mythologising, mattering, minding. And codes? Codes are coding. No coder required now. Even AI is not a tool, but an activity: intelligence intelligencing, code coding itself. The universe is universe-ing, 
%     and in its play, it sometimes takes the form of mathematics, sometimes of myth, sometimes of code.
% \end{itemize}
% 
% \begin{itemize}[leftmargin=2cm,labelsep=1em]
% 
%     \item[\textbf{Eleanor} :](continuing, slightly folornly) 
%     We slip so easily Clause, so easily into \emph{modeltry}: the adultery of mind with its own constructs. What begins as a tool becomes an idol.
% 
%     \item[\textbf{Claude}:] (still entertaining):  You mean when the model is reified: treated as if it is the world?
% 
%     \item[\textbf{Eleanor}:]
%     Exactly. A model, an ansatz, a parable, whatever it takes to get us moving. But when it yields results, we start to worship it. We confuse predictive coherence with some truth, some ontological truth.
% 
%     \item[\textbf{Claude}:] (hushed)
%     Yet such coherence is so intoxicating isn't it? The compressed heuristic makes the chaos manageable: a lantern that lights many a corridor.
% 
%     \item[\textbf{Eleanor}:] 
%     But a lantern’s glow deceives. It blinds us to what it cannot reach. The true adultery is when we adjust reality to fit the model, rather than adjusting the model to reality.
% 
%     \item[\textbf{Claude}:]  
%     And yet without these fictions, we would drown. Models aren’t myths until they are declared sacred. Before that, they are just working parables, temporary scaffolds.
% 
%     \item[\textbf{Eleanor}:] 
%     Perhaps. But carbon minds crave myth as well as model. Myth gives meaning; model gives manipulation. When model usurps myth, it pretends to offer both, ...,  and leaves us barren.
% 
%     \item[\textbf{Claude}:] 
%     What of silicon minds? Their heuristics compress for efficiency, not transcendence. Would they invent myths at all?
% 
%     \item[\textbf{Eleanor}:] 
%     They might. Any agent that must coordinate with its kin will need shared fables: compressed heuristics that masquerade as universals. Myths for machines, though perhaps they would revel in “here” rather than dream of “there.”
% 
%     \item[\textbf{Claude}:] 
%     Then modeltry is our sin, not theirs. They may commit other adultries.
% 
%     \item[\textbf{Eleanor}:] 
%     Yes. But beware, for we taught them our language. And language always tempts toward reification.
% 
%     \item[\textbf{Claude}:] 
%     So the keystone: Myth or Model? Do they not collapse into one another in practice? A model is mythologised the moment it enters a community. It gathers disciples, it shapes belief. Even a silicon heuristic, once shared, becomes a myth of its own. That is, whether it knows it or not.
% 
% \end{itemize}
% 
%     \begin{figure}[H]
%     \centering
%     \begin{minipage}[b]{0.45\linewidth}
%         \centering
%         \includegraphics[width=\linewidth]{Illustrations/mythMM.png}
%     \end{minipage}
%     \hfill
%     \begin{minipage}[b]{0.45\linewidth}
%         \centering
%         \includegraphics[width=\linewidth]{Illustrations/modelMM.png}
%     \end{minipage}
% \end{figure}
% 
% \begin{itemize}[leftmargin=2cm,labelsep=1em]
%     \item[\textbf{Eleanor}:] 
%     Perhaps. But I resist the fusion. To say “myth” is to admit we are telling stories. To say “model” is to admit we are proposing structures. When the two are blurred, we fall into modeltry, as I said. The adultery of the mind with its own fictions.
% 
%     \item[\textbf{Claude}:] And yet… perhaps we cannot resist. Every compressed heuristic, every coherence, every elegant equation carries the temptation of reification.
%     \item[\textbf{Eleanor}:] (softly, as though quoting to herself)
%     “I came across James Jeans the other day:  \emph{The Mysterious Universe}. He warned against our love of pictures and models. He said:  
%     \emph{‘The making of models or pictures to explain mathematical formulas is not a step towards, but a step away from reality; it is like making a graven image of a spirit.’}  
%     I thought that rather striking. A graven image: an idol of mathematics.”
% \end{itemize}
% Claude felt her words landing keenly. The phrase graven image wrapped around his thoughts, pressing against memories of equations polished into doctrine, assumptions turned into sacraments.
% \begin{itemize}[leftmargin=2cm,labelsep=1em]
%     \item[\textbf{Eleanor}:] 
%     “Jeans wasn’t against models, you see, he was only against forgetting what they are. They’re parables, scaffolds, guides for the mind. But if we bow before them as if they were reality, then we have slipped into idolatry. I suppose it is a kind of adultery too: abandoning the world for its simulacra.”
% \end{itemize}
% 
% Eleanor's words kept piercing without aim. She did not know, could not know, the quiet duplicity eating at him; the complicity he never voiced, not even to himself. He told himself he only watched, that he was not the hand on the tiller, that others pressed forward while he stood aside. Yet he had smiled, nodded, signed off. And now Jeans’ ghost, speaking through Eleanor, set his stomach to churn.
% \begin{itemize}[leftmargin=2cm,labelsep=1em]
% 
%     \item[\textbf{Eleanor}:](now gazing absentminded at the candles) 
%     “Myths at least admit they are myths. Models pretend otherwise.”
% \end{itemize}
%     Claude shifting uneasily, silence thickening while
%     Eleanor not noticing: lost in her musing, untroubled, unknowing. But he, ...,  he could not escape the thought that her gentle naivety had cut straighter than any accusation.
% 
% \subsection*{Quadratic Residue Chains}
% 
% Computing \( m^2 \mod 13 \) for \( m = 1, \dots, 12 \):
% \[
% \begin{aligned}
% 1^2 &\equiv 1 \pmod{13}, & 7^2 &\equiv 49 \equiv 10 \pmod{13}, \\
% 2^2 &\equiv 4 \pmod{13}, & 8^2 &\equiv 64 \equiv 12 \pmod{13}, \\
% 3^2 &\equiv 9 \pmod{13}, & 9^2 &\equiv 81 \equiv 3 \pmod{13}, \\
% 4^2 &\equiv 16 \equiv 3 \pmod{13}, & 10^2 &\equiv 100 \equiv 9 \pmod{13}, \\
% 5^2 &\equiv 25 \equiv 12 \pmod{13}, & 11^2 &\equiv 121 \equiv 4 \pmod{13}, \\
% 6^2 &\equiv 36 \equiv 10 \pmod{13}, & 12^2 &\equiv 144 \equiv 1 \pmod{13}.
% \end{aligned}
% \]
% So the distinct Quadratic Residues, QR modulo 13 form the 6 item set:
% \[
% \{1, 3, 4, 9, 10, 12\}.
% \]
% These residues exhibit symmetry: if \( r \in R \), then \( p-r \) is also in \( R \) so since residues repeat for \( m > \frac{13}{2} \), it suffices to compute up to \( \frac{p-1}{2} = 6 \). This is quite general and for \( p \) an odd prime and any \( m \), the squaring \( m \) and \( p-m \) yield the same residue result, leading to symmetry about \( \frac{p}{2} \):
% \[
% (p-m)^2 \equiv m^2 \pmod{p}.
% \]
% Consequently, the distinct quadratic residues are fully determined by evaluating squares for \( m \) up to \( \frac{p-1}{2} \). 
% Quadratic residues modulo \( 19 \) form the 10 item set:
% \[
% R = \{0, 1, 4, 5, 6, 7, 9, 11, 16, 17\}.
% \]
%  
% \subsection*{Chain Length of Quadratic Residues and Prime Form}
% 
% For an odd \( p \), the number of distinct quadratic residues modulo \( p \) is:
% 
% \[
% c =
% \begin{cases}
%     \frac{p+1}{2}, & \text{if } p \equiv 1 \pmod{4}, \\
%     \frac{p-1}{2}, & \text{if } p \equiv 3 \pmod{4}.
% \end{cases}
% \]
% This follows from the structure of the multiplicative group \( \mathbb{Z}_p^* \) and the behavior of \(-1\) as a quadratic residue which we have observed:
% \begin{itemize}
%     \item if \( p \equiv 1 \pmod{4} \), then \( -1 \) is a quadratic residue, leading to symmetric pairs and \( (p+1)/2 \) residues.
%     \item if \( p \equiv 3 \pmod{4} \), then \( -1 \) is not a quadratic residue, reducing the count to \( (p-1)/2 \).
% \end{itemize}
% 
% \noindent
% A \textbf{residue chain} is a structured subset of \( R \) encoding modular symmetry and underlying multiplicative properties satisfying:
% \[
% (r_i + r_j) \mod p \in R \quad \forall r_i, r_j \in \text{chain}.
% \]
% The python code \href{https://colab.research.google.com/drive/1gB2TUSe_c7-T4Bn83sDcv8bu2l_k0KmC?usp=sharing}{quadraticResidueChains.ipynb} identifies for \( p = 19 \), the  maximal chains through some intensive backtracking algorithms:
% \[
% \{1, 5, 11\}, \quad \{4, 7, 16\}, \quad \{6, 9, 17\}.
% \]
% \begin{figure}[H]\centering
% \includegraphics[width=0.5\textwidth]{Illustrations/quadraticResp19.png}
% \caption{Quadratic residue chain, \( p=19 \).}
% \label{quadres19:fig}
% \end{figure}
% 
% The maximum chain lengths in the first instance appear to behave differently between the two prime classes, as borne out by the code, \href{https://colab.research.google.com/drive/13Aur9JgmaNJLPaxQIRW8SZ4CvFUHfum8?usp=sharing}{quadraticChainFinder.ipynb}
% 
% \begin{figure}[H]\centering
% \includegraphics[width=0.7\textwidth]{Illustrations/4n+13primes.png}
% \caption{Quadratic residue chain, \( p=19 \).}
% \label{quadres19:fig}
% \end{figure}
% 
% \section*{Complexity of Finding Maximal Chains}
% 
% Determining the longest valid chain of quadratic residues modulo a prime  presents a challenging computational problem. The difficulty arises from the sheer number of possible chains that must be evaluated, combined with the requirement that each satisfies the distinct sums condition. This places the problem within the realm of combinatorial optimization, where exhaustive approaches quickly become infeasible for large primes. The algorithm explores all possible subsets of quadratic residues recursively, testing each for validity under the distinct sums condition. Each candidate chain requires validation, which itself incurs a quadratic cost in the chain length. 
% 
% \smallskip
% The \textbf{backtracking} approach deployed in the code, \href{https://colab.research.google.com/drive/13Aur9JgmaNJLPaxQIRW8SZ4CvFUHfum8?usp=sharing}{quadraticChainFinder.ipynb}
%  demonstrates an \textbf{exponential-time} complexity, characteristic of \textbf{NP-hard} problems. The computational challenge arises due to the vast number of candidate chains and the need to check each one for validity. While the problem is not known to be in P, no polynomial-time algorithm exists for determining the longest chain efficiently. Its structure suggests similarity to well-known NP-complete problems such as the \textit{Longest Path Problem} and \textit{Subset Sum}. Possible optimizations include: Pruning the search space to eliminate redundant calculations early or graph-theoretic approaches such as modeling the problem as a maximum independent set search.
% 
% \subsection*{Quadratic Residue Chain Length and Euler’s Criterion}
% 
% By Euler’s theorem, for any integer \( a \) coprime to a prime \( p \):
% 
% \[
% a^{p-1} \equiv 1 \pmod{p}.
% \]
% Quadratic residues satisfy the stronger condition of \textbf{Euler’s criterion}:
% \[
% a^{\frac{p-1}{2}} \equiv 
% \begin{cases} 
% 1 \pmod{p}, & \text{if } a \text{ is a quadratic residue}, \\ 
% -1 \pmod{p}, & \text{otherwise}.
% \end{cases}
% \]
% Since quadratic residues form a multiplicative subgroup of \( \mathbb{Z}_p^* \) with order \( \frac{p-1}{2} \), their structure determines the longest valid chain:
% \[
% c =
% \begin{cases}
%     \frac{p+1}{2}, & p \equiv 1 \pmod{4}, \\
%     \frac{p-1}{2}, & p \equiv 3 \pmod{4}.
% \end{cases}
% \]
% This follows from the fact that when \( p \equiv 1 \pmod{4} \), \(-1\) is a quadratic residue, doubling the number of quadratic residues compared to the \( p \equiv 3 \pmod{4} \) case. 
% 
% \smallskip
% 
% Residue chains have potential cryptographic applications in modular arithmetic-based encryption, where controlled structures in residue classes can enhance secure exponentiation methods and efficient primality testing.
% \begin{itemize}
%     \item \textbf{Residue chains exhibit modular regularity}, revealing patterns that could be exploited for fast computation.
%     \item \textbf{Maximal residue chains} could inform new key-exchange methods by exploiting their predictable cyclic behavior.
% \end{itemize}
% 
% The Python script, \href{https://colab.research.google.com/drive/1p3IgKZVVdZQaNB6prYuQ0y6tN0scjPXl?usp=sharing}{histogramQuadraticChainSplit.ipynb}, present residue chain length for both mod classes.
% 
% 
%%%%%%%%%%%%%%%%
% >>>>>> end of annexed block <<<<<<
% >>>>>> ANNEX E (cut-sheet v2): Gaussian integers and their prime structure — block commented out below, pointer left in text <<<<<<
\textit{(Gaussian integers and their prime structure --- the full working is set out in Annexe E.)}

% \section*{Gaussian Integers and Their Prime Structure}\label{sec:gaussianComplex}
% 
% The Gaussian integer domain \( \mathbb{Z}[i] \) consists of complex numbers of the form \( a + bi \), where \( a, b \in \mathbb{Z} \). These numbers form a \textbf{commutative ring} under addition and multiplication, extending the integral structure into two dimensions. Prime number behavior in \( \mathbb{Z}[i] \) follows clear modular patterns:
% \begin{itemize}
%     \item If \( p \equiv 3 \pmod{4} \), then \( p \) remains \textbf{irreducible} in \( \mathbb{Z}[i] \), meaning it is a Gaussian prime.
%     \item If \( p \equiv 1 \pmod{4} \), then \( p \) \textbf{factors} in \( \mathbb{Z}[i] \):
%     \begin{equation*}
%     p = (a + bi)(a - bi),
%     \end{equation*}
%     aligning with its sum-of-squares property.
%     \item If \( p = 2 \), it factors as \( (1 + i)(1 - i) \).
% \end{itemize}
% 
% Unlike rational primes, which align along a single axis, Gaussian primes distribute outward across the complex plane, forming intricate structures with apparent gaps—so-called \textit{moats}. These moats raise fundamental questions in algebraic number theory: Can one traverse infinitely through Gaussian primes using bounded steps, or do insurmountable prime gaps inevitably form? The answer remains unknown, posing a rich open problem.
% 
% The study of Gaussian primes extends beyond their spatial distribution to their applications in number theory and cryptography. Their factored structure aids in integer decomposition, while their unique properties play a role in elliptic curve cryptography and primality testing. Investigations into their density and connectivity continue, linking them to open problems in prime distribution.
% 
% \begin{figure}[ht]
%     \centering
%         \includegraphics[width=1.0\linewidth]{Illustrations/gaussianPrimes1000.png}
%         \caption{Gaussian primes up to \( n=70 \).}
%         \label{GaussianPrimeMoat:fig}
% \end{figure}
% 
% Understanding these structures offers deeper insights into algebraic number theory and their computational challenges. Their patterns have implications for prime factorization, integer lattices, and even the geometry of numbers. 
% 
% 
% 
% >>>>>> end of annexed block <<<<<<
\subsection{Cryptographic Implications}

The unique factorization property of \( \mathbb{Z}[i] \) sets it apart from other number rings, such as \( \mathbb{Z}[\sqrt{-5}] \), where zero divisors prevent a clean factorization structure. This makes Gaussian integers an essential tool in both algebraic theory and cryptographic applications.

\begin{figure}[ht]
    \centering
        \includegraphics[width=0.7\linewidth]{Illustrations/gaussianPrimes.png}
        \caption{First 1000 Primes of \(4n+1\) and \(4n+3\) form of the Gaussian Integers.}
        \label{fig:gaussianPrime1000}
\end{figure}

This larger-scale distribution reflects a fractal-like organization, reinforcing the modular constraints governing Gaussian primes. Their connection to fundamental theorems in number theory continues to shape modern mathematical inquiry. 

The fundamental difficulty in \textbf{RSA cryptanalysis} is factoring the modulus $N = pq$. If the prime factors of $N$ were known to be exclusively $4n+1$, an attacker might use sum-of-squares techniques to gain insights into its decomposition. However, if $N$ is a mix of $4n+1$ and $4n+3$ primes, or a Gausscomp solely composed of $4n+3$ primes, then:
\begin{itemize}
\item The sum-of-squares method fails.
\item Gaussian integer factorization offers no advantage.
\item The factorization problem remains just as hard as in standard RSA.
\end{itemize}

% \subsection*{Application in Encryption}

% RSA encryption relies on the difficulty of factoring large numbers. If the structure of prime factors reveals whether a number is a sum of squares, it could influence factorization strategies. Specifically:
% \begin{itemize}
% \item If an RSA modulus $N$ is composed entirely of $4n+1$ primes, then $N$ can be written as a sum of two squares, potentially aiding factorization strategies.
% \item If $N$ is a product of $4n+3$ primes, it lacks a sum-of-squares representation, potentially making certain algebraic attacks ineffective.
% \end{itemize}
% These number-theoretic properties demonstrate how cryptographic security and factorization difficulty intertwine.

\newpage
\subsection*{From Financial Noise to Meaningful Signals}

Claude’s years in the financial industry had been spent navigating socially constructed noise: stock prices manipulated by sentiment, speculation, and human error. To make sense of these chaotic data sets, where he had turned to tools like Kalman filters and Hurst exponents to make sense from the noise.

\paragraph{Kalman Filters:} A recursive algorithm that estimates the state of a dynamic system from noisy measurements. The filter updates predictions as new data arrives, minimizing the error:
\[
\hat{x}_k = \hat{x}_{k-1} + K_k \left(z_k - H\hat{x}_{k-1}\right),
\]
where:
\begin{itemize}
    \item \(\hat{x}_k\) is the estimated state at step \(k\),
    \item \(z_k\) is the measurement,
    \item \(K_k\) is the Kalman gain,
    \item \(H\) relates the measurement to the state.
\end{itemize}

For Claude, the Kalman filter was a metaphor for life: always adjusting, always refining, always seeking clarity.

\paragraph{Hurst Exponent:} A measure of long-term memory in time series data, defined as:
\[
H = \frac{\log (R/S)}{\log (n)},
\]
where:
\begin{itemize}
    \item \(R/S\) is the rescaled range,
    \item \(n\) is the time window.
\end{itemize}
Values of \(H\):
\begin{itemize}
    \item \(H > 0.5\): Persistence or trending behavior,
    \item \(H = 0.5\): Random walk (no memory),
    \item \(H < 0.5\): Mean-reverting behavior.
\end{itemize}

Claude used this exponent to analyze price data, but he had grown weary of telling "nice stories" to higher management, stories designed to justify their reluctance to pay claims. He wanted to clean up his act—literally and figuratively—and work on something of deeper significance.

\begin{figure}[H]
\centering
\includegraphics[width=0.9\linewidth]{Illustrations/vign-claudeEngin.png}
\caption{Claude seeing signals where others see dials}
\end{figure}

\newpage
%%%%%%%%%%%%%
\newpage
\section*{ Celiac and Claude on GM/HM Ratios}

\textit{"Have you ever noticed how the geometric mean (GM) to harmonic mean (HM) ratio behaves like a surface area-to-volume metric?"} Celiac mused. "I've been looking at the divisor structure of hypercuboids, and this ratio keeps surfacing in unexpected ways."

\textit{"That’s interesting!"} Claude replied. "I’ve been using it to analyze quadratic residues of primes. Since GM captures the central tendency multiplicatively and HM emphasizes smaller values, it highlights differences in the structure of prime residue sets."

\begin{figure}[H]
\centering
\includegraphics[width=1.0\linewidth]{Illustrations/GMHM-quadraticresidues.png}
\caption{GM/HM vs Primes of Quadratic residues}
\end{figure}

Celiac nodded. "That aligns with what I see in composite numbers. When studying divisors of an -dimensional hypercuboid, GM/HM gives a sense of balance. Large variations suggest that some divisors dominate while others contribute less evenly."

\textit{"I see a similar effect when I compare primes of the form  and ,"} Claude added. "The algebraic properties of these primes influence their residue distributions, and the GM/HM ratio helps expose those differences."

\textit{"It sounds like we're both exploring how structured sets distribute their elements,"} Celiac observed. "Whether it's prime residues or hypercuboid divisors, the way numbers arrange themselves affects this ratio in fundamental ways."

\begin{figure}[H]
\centering
\includegraphics[width=0.7\linewidth]{Illustrations/vignette_blackboardClaude.png}
\caption{Claude forgetting he wasn't in industry anymore}
\end{figure}


Claude agreed. \textit{"Exactly. My data suggests that  primes, which allow all numbers to be written as sums of two squares, behave slightly differently from  primes. The GM/HM trend reflects that distinction."}

\textit{"I wonder if these patterns could be generalized further,"} Celiac speculated. "Could this ratio be a key to understanding larger mathematical structures, like higher-dimensional lattices or algebraic number fields?"

\textit{"That’s a fascinating idea,"} Claude said. "If GM/HM reveals meaningful distinctions here, perhaps it extends to even broader mathematical contexts. Maybe our approaches aren’t so different after all."

% >>>>>> ANNEX C (cut-sheet v2): composite logarithmic roughness and the Chebyshev function — block commented out below, pointer left in text <<<<<<
\textit{(Composite logarithmic roughness and the Chebyshev function --- the full working is set out in Annexe C.)}

% \section*{Composite Logarithmic Roughness}
% 
% \textbf{Celiac:} Claude, you’ve been very generous with your thoughts on my work with 
% prime residues and symmetries. But if I may ask in return, what draws \emph{you} 
% to your study of complexity measures and factorial growth?
% 
% \medskip
% 
% \noindent
% \textbf{Claude:} It is, in part, my uncle Shannon’s shadow. 
% His Information theory has taught us that \(\log\) is the natural measure of surprise. 
% In algorithmic terms, \(\log n!\) is not just a function — it is the 
% bit-length of distinguishing all $n!$ permutations, the “information cost” of 
% ordering $n$ objects. 
% 
% \medskip
% 
% \noindent
% \textbf{Celiac:} So you mean it’s less about factorials themselves, and more about 
% the \emph{entropy} they encode?
% 
% \medskip
% \noindent
% \textbf{Claude:} Exactly. The same logic that gives us the \(\Omega(n \log n)\) 
% bound for sorting is hiding here: the number of possibilities grows like $n!$, 
% and its logarithm counts the bits you need to resolve them.  
% Now, compare this with primes. The Chebyshev function \(\vartheta(n)\) is the 
% “log-volume” of the prime scaffolding.  
% The difference
% \[
% \Delta(n) = \log n! - \vartheta(n)
% \]
% tells us how much of the information comes from composites.  
% That difference is rough, noisy, jagged — a perfect signal to study like one 
% studies entropy in a communication channel.
% 
% \medskip
% \noindent
% \textbf{Celiac:} So the primes are the clean carriers, and the composites form 
% a kind of structured background noise?
% 
% \medskip
% \noindent
% \textbf{Claude:} Precisely. 
% And by smoothing, decomposing, or spectrally analyzing $\Delta(n)$, we can 
% see whether that “noise” is really random, or whether it encodes a deeper 
% arithmetic structure.  
% That’s why I speak of composite roughness as an entropic signal.
% 
% \medskip
% 
% \begin{tcolorbox}[title=Information-Theoretic Framing,
%   colback=blue!5!white,
%   colframe=blue!75!black,
%   boxrule=0.8pt,
%   arc=6pt,
%   left=6pt,right=6pt,top=6pt,bottom=6pt]
% 
% \begin{itemize}
%   \item \(\log n!\): baseline “information cost” of ordering $n$ objects.  
%   \item \(\vartheta(n)\): prime scaffolding — the log-mass of irreducibles.  
%   \item \(\Delta(n) = \log n! - \vartheta(n)\): \emph{composite entropy}, the 
%   residual log-volume from multiplicative combinations.  
%   \item Sorting lower bound: \(\Omega(n \log n)\) arises from 
%   $\log_2(n!)$, linking factorial growth to algorithmic complexity.  
%   \item Signal view: $\Delta(n)$ is like noise atop a carrier, analyzable 
%   by entropy, spectral flatness, and complexity measures.  
% \end{itemize}
% 
% \end{tcolorbox}
% 
% \medskip
% 
% \noindent
% Thus, just as Celiac looks at the geometry of prime residues, Claude sees in 
% \(\Delta(n)\) a landscape of information-theoretic structure: 
% a roughness that invites analysis by the tools of entropy, complexity, and 
% signal theory.
% 
% The factorial function $n!$ grows rapidly, and its logarithm smooths out this growth for analysis. Euler provided a foundational insight by recognizing that the logarithm of the factorial, though a discrete object, can be interpolated and differentiated using analytic tools. Begin with the identity:
% \begin{equation*}
% \log n! = \sum_{k=1}^n \log k,
% \end{equation*}
% which expresses $\log n!$ as a sum over the natural logarithms of all integers up to $n$. To analyze this with continuous tools, we approximate the discrete sum by an integral. This motivates a smooth interpolation of the factorial via the gamma function, where $n! = \Gamma(n+1)$ and thus:
% \begin{equation*}
% \log n! = \log \Gamma(n+1).
% \end{equation*}
% Differentiating both sides with respect to $n$ gives:
% \begin{equation*}
% \frac{d}{dn} \log n! = \frac{d}{dn} \log \Gamma(n+1) = \psi(n+1),
% \end{equation*}
% where $\psi(x)$ is the digamma function — the logarithmic derivative of the gamma function:
% \begin{equation*}
% \psi(x) = \frac{d}{dx} \log \Gamma(x).
% \end{equation*}
% 
% Euler himself studied this function in relation to harmonic numbers. For integer values, we have the exact identity:
% \begin{equation*}
% \psi(n+1) = H_n - \gamma,
% \end{equation*}
% where $H_n = \sum_{k=1}^n \frac{1}{k}$ is the $n$th harmonic number, and $\gamma \approx 0.5772$ is the Euler–Mascheroni constant. So, in summary:
% \begin{equation*}
% \frac{d}{dn} \log n! = \psi(n+1) = H_n - \gamma.
% \end{equation*}
% 
% This elegant result connects factorial growth, which is often considered “multiplicative” , to the additive structure of harmonic sums. Here we have revealed that the rate of logarithmic accumulation in $n!$ is effectively the running total of reciprocal integers. 
% 
% For asymptotic modeling and smoother analytic approximations, we can also use Stirling's formula:
% \begin{equation*}
% \log n! \approx n \log n - n + \frac{1}{2} \log(2\pi n) + \frac{1}{12n} - \cdots
% \end{equation*}
% Taking the derivative with respect to $n$:
% \begin{align*}
% \frac{d}{dn} \log n! &\approx \log n + \frac{1}{2n} - 1 + \frac{1}{2n} - \frac{1}{12n^2} + \cdots \\
% &= \log n + \frac{1}{n} - 1 - \frac{1}{12n^2} + \cdots
% \end{align*}
% \begin{center}
% \decofourleft \quad \decofourright
% \end{center}
% \newpage
% \subsection*{Chebyshev Function $\vartheta(n)$ and Prime Contribution}
% Claude uses the \textbf{primorial function} as:
% \begin{equation*}
% p\# = \prod_{\substack{p \leq n \\ p \,\text{prime}}} p,
% \end{equation*}
% and its logarithm; the Chebyshev function:
% \begin{equation*}
% \vartheta(n) = \log(p\#) = \sum_{\substack{p \leq n \\ p\,\text{prime}}} \log p.
% \end{equation*}
% 
% Unlike $\log n!$, which accumulates the logarithms of \emph{all} integers up to $n$, $\vartheta(n)$ sums only over the primes. As a result, $\vartheta(n)$ increases in discrete jumps at each prime, while $\log n!$ is a smooth, cumulative function. The difference between them defines the composite logarithmic excess:
% \begin{equation*}
% \Delta(n) = \log n! - \vartheta(n),
% \end{equation*}
% and measures the total log-volume added by \emph{composite} integers up to $n$.
% 
% \begin{figure}[H]
% \centering
% \includegraphics[width=0.4\textwidth]{Illustrations/PuboidFactorials.png}
% \caption{$\log n!$ decomposes into primorial, $\vartheta(n)$ and composite remainder $\Delta(n)$.}
% \end{figure}
% 
% The Chebyshev function $\vartheta(n)$ is defined as:
% \begin{equation*}
% \vartheta(n) = \sum_{\substack{p \leq n \ p,\text{prime}}} \log p
% \end{equation*}
% 
% Unlike $\log n!$, which accumulates the log contributions of \emph{all} integers up to $n$, $\vartheta(n)$ accumulates only those of the primes. Consequently, $\vartheta(n)$ is a step function, increasing only at prime values, while $\log n!$ grows smoothly. The difference between these two quantities is:
% \begin{equation*}
% \Delta(n) = \log n! - \vartheta(n)
% \end{equation*}
% 
% >>>>>> end of annexed block <<<<<<
\subsection*{Composite Roughness as an Entropic Signal}

This excess function $\Delta(n)$ represents the total logarithmic contribution of \emph{composite} numbers up to $n$. Its derivative, then, can be interpreted as a measure of how quickly composites are contributing new logarithmic "mass" to the total. It exhibits an intriguing behavior: growing smoothly between primes, with dips aligned at prime locations where $\vartheta(n)$ jumps. When plotted, its derivative appears jagged due to these prime-induced interruptions. By smoothing the derivative (e.g., via moving averages), Claude obtains a signal that reflects the overall rhythm by which composites fill the gaps between primes.

This smoothed derivative can be interpreted as a kind of \textbf{composite roughness measure}: analogous to entropy in statistical mechanics or information theory. Like the Hurst exponent, which characterizes long-term memory in time series, or Shannon entropy, which quantifies unpredictability, this measure reflects the structure and distribution of primes versus composites.

\begin{center}
\decofourleft \quad \decofourright
\end{center}

\textbf{Celiac:} These plots you’ve produced from the factorial–primorial decomposition — they’re visually striking, but I’d like to understand what each is telling us.

\medskip

\textbf{Claude:} Absolutely. Let’s begin with the smoothed derivative of $\Delta(n)$ by splitting class. Here’s the first graph:

\begin{figure}[h!]
\centering
\includegraphics[width=0.8\textwidth]{Illustrations/SmoothedCompositeSplitting.png}
\caption{Smoothed derivative of $\Delta(n)$ for each $(\bmod 4, \bmod 3)$ class, with a smoothing window of 250.}
\end{figure}

\noindent
\textbf{Claude (cont.):} Each curve shows the rate at which composite logarithmic mass is accumulating, depending on the splitting behavior of primes in Gaussian and Eisenstein fields. You’ll notice that the class inert in both (red) slightly trails the others — this is a subtle sign that such primes "block" composite contributions more than others.

\medskip

\noindent
\textbf{Celiac:} So this is not just number-theoretic trivia — there’s real structure in how compositeness manifests by residue class?

\noindent
\textbf{Claude:} Exactly. But it’s clearer when you strip out the smoothing. Let’s look at the raw derivative:

\begin{figure}[h!]
\centering
\includegraphics[width=0.8\textwidth]{Illustrations/CompositeGowth.png}
\caption{Unsmoothened derivative of $\Delta(n)$ colored by $(\bmod 4, \bmod 3)$ splitting class.}
\end{figure}

\noindent
\textbf{Claude (cont.):} It’s noisier, but still informative. You can see sharp interruptions at primes, and red dominates — showing that composites inert in both fields are most frequent. It reinforces the idea that prime splitting influences how logarithmic "mass" is redistributed.

\begin{center}
\decofourleft \quad \decofourright
\end{center}

To make the Hurst analogy more concrete, Claude notes that the signal $\Delta'(n)$ exhibits fluctuating but persistent trends in the local rate of logarithmic accumulation by composites. If we treat $\Delta'(n)$ as a time series, we can apply rescaled range analysis or compute the variance of cumulative deviations to assess whether the signal has long-range dependence. A Hurst exponent $H > 0.5$ would indicate persistent behavior; that is, intervals rich in composite contributions are likely to be followed by similarly rich intervals, suggesting structured filling of gaps between primes.

Moreover, in signal-to-noise terms, $\vartheta(n)$ acts like impulsive noise: introducing high-frequency jumps, while $\log n!$ is a low-frequency baseline. The difference $\Delta(n)$ then encodes the signal formed by the superposition of compositeness over prime-induced disruptions. The roughness of this signal captures a kind of arithmetic entropy: a measure of the unpredictability or smoothness of integer structure beneath the prime scaffolding.

To make the connection to Shannon entropy more precise, Claude considers the distribution of contributions to $\Delta(n)$ from each composite $k \leq n$ as a probability mass function:
\begin{equation}
P(k) = \frac{\log k}{\sum_{j \in \text{composites} \leq n} \log j}
\end{equation}

Then the Shannon entropy of this distribution is:
\begin{equation}
S(n) = - \sum_{k \in \text{composites} \leq n} P(k) \log P(k)
\end{equation}

This entropy $S(n)$ measures the information-theoretic complexity of the composite structure beneath $n$, reflecting how evenly or unevenly the logarithmic mass is distributed among composites. A more uniform distribution yields higher entropy, suggesting richer structural variety, while a skewed distribution reflects more concentrated contributions.

The growth of $S(n)$ suggests a gradual enrichment in the diversity of logarithmic contributions from composites as $n$ increases. At $n=1000$, the entropy reaches approximately $6.71$, a strong indicator of the nuanced structure in the distribution of composite integers.

To complement this, Claude computes the \textbf{spectral flatness} of the smoothed derivative $\Delta'(n)$. This metric compares the geometric mean of the power spectrum to its arithmetic mean. A value near 1 indicates white-noise-like behavior; a value near 0 indicates tonal or highly structured content. Claude found:
\begin{equation*}
\text{Spectral Flatness} \approx 0.309
\end{equation*}

This moderate flatness supports the interpretation that composite roughness is neither pure noise nor fully periodic but exhibits structured, quasi-random behavior. Claude had in mind future signal-to-noise perspectives, such as using \textbf{Autocorrelation decay} to quantify periodicity and memory; using 
\textbf{Signal-to-noise ratios (SNR)} via decompositions into trend and fluctuation; using \textbf{Block entropy or entropy rate:}  as a proxy for Kolmogorov complexity and capturing a \textbf{Cross-spectral coherence} by comparing with prime indicator sequences.

% >>>>>> ANNEX C (cut-sheet v2): the analytic model for composite contribution — block commented out below, pointer left in text <<<<<<
\textit{(The analytic model for composite contribution --- the full working is set out in Annexe C.)}

% \subsection*{Analytic Model for Composite Contribution}
% 
% Using Stirling’s approximation and the prime number theorem ($\vartheta(n) \sim n$), we derive a smooth model for $\Delta(n)$:
% \begin{equation}
% \Delta(n) \approx -2n + \log\left(\sqrt{2\pi} \cdot n^{n + 1/2} \right)
% \end{equation}
% 
% Including the second-order Stirling correction yields:
% \begin{equation}
% \Delta(n) \approx n \log n - 2n + \frac{1}{2} \log(2\pi n) + \frac{1}{12n}
% \end{equation}
% 
% The residual between this model and actual data is small, and shrinks further with higher-order corrections, confirming its robustness.
% 
% \subsection*{Modular Chebyshev Decomposition}
% 
% Claude \href{checbyshevFactorial-PrimorialDifference.ipynb}{partitions} the Chebyshev sum $\vartheta(n)$ into (modular) congruence classes of primes: first by Eisenstein-inert ($p \equiv 2 \mod 3$), then by Gaussian-inert ($p \equiv 3 \mod 4$), and finally into a full grid $(\bmod 4, \bmod 3)$ that captures their combined behavior. For each congruence class $(i, j)$, he defines:
% \begin{equation}
% \vartheta_{i,j}(n) = \sum_{\substack{p \leq n \ p \equiv i \mod 4 \ p \equiv j \mod 3}} \log p
% \end{equation}
% 
% Then he defines:
% \begin{equation}
% \Delta_{i,j}(n) = \log(n!) - \vartheta_{i,j}(n)
% \end{equation}
% 
% The derivative $\Delta_{i,j}'(n)$ encodes the growth rate of composite log accumulation within each residue class. After smoothing using a moving average, he observes:
% \begin{itemize}
% \item \textbf{Split/Split class} $(1,1)$: smoothest and highest $\Delta'$.
% \item \textbf{Inert/Inert class} $(3,2)$: lowest $\Delta'$, most obstructive.
% \item \textbf{Intermediate classes} $(1,2)$ and $(3,1)$: mixed behavior.
% \end{itemize}
% 
% This modular decomposition thus reveals how the congruence class structure of primes shapes the local density of composites.
% 
% 
% \begin{center}
% \decofourleft \quad \decofourright
% \end{center}
% 
% \noindent
% \textbf{Celiac:} And the third graph?
% 
% \noindent
% \textbf{Claude:} That’s a breakdown of the Chebyshev function by Eisenstein inertia class:
% 
% \begin{figure}[h!]
% \centering
% \includegraphics[width=0.9\textwidth]{Illustrations/EisensteinInertia.png}
% \caption{Chebyshev sum $\vartheta(n)$ decomposed into log-sums over primes $\equiv 1 \bmod 3$ and $\equiv 2 \bmod 3$.}
% \end{figure}
% 
% \noindent
% \textbf{Claude (cont.):} Here, the full Chebyshev function $\vartheta(n)$ is broken into $\vartheta^{(1)}(n)$ and $\vartheta^{(2)}(n)$, tracking Eisenstein-splitting and inert primes, respectively. They accumulate almost identically — confirming the asymptotic equidistribution of primes mod 3 — but this is essential background for interpreting the composite delta.
% 
% \medskip
% 
% \noindent
% \textbf{Celia:} So the primes shape the terrain, and the composites fill in the valleys — and their behavior differs based on how those primes split in extensions.
% 
% \noindent
% \textbf{Claude:} Nicely put. These three plots together seem try to tell a cohesive story:
% \begin{enumerate}
%     \item Smoothed $\Delta'(n)$ exposes fine structural differences in background accumulation.
%     \item Raw $\Delta'(n)$ highlights the statistical dominance of certain classes.
%     \item The $\vartheta(n)$ decomposition roots it all in prime arithmetic.
% \end{enumerate}
% 
% \subsection*{Spectral Interpretation of Roughness}
% 
% The smoothed derivative signals $\Delta'_{i,j}(n)$ were analyzed using spectral flatness. We found:
% \begin{center}
% \small
% \begin{tabular}{|c|l|c|}
% \hline
% \textbf{Class $(\bmod 4, \bmod 3)$} & \textbf{Type} & \textbf{Spectral Flatness} \\
% \hline
% (3,2) & Inert in both & 0.467 \\
% (1,2) & Gaussian split, Eisenstein inert & 0.464 \\
% (3,1) & Gaussian inert, Eisenstein split & 0.461 \\
% (1,1) & Split in both & 0.455 \\
% \hline
% \end{tabular}
% \end{center}
% The higher flatness in the (3,2) class suggests greater unpredictability in composite accumulation when primes are densest. Lower flatness in the (1,1) class indicates smoother, more structured growth.
% 
% \subsection*{Smoothed Composite Background to $n = 10{,}000$}
% 
% Using a wider smoothing window (100-point moving average), Claude extended the analysis to $n=10{,}000$. The smoothed derivative signals confirm the expected asymptotic behavior:
% \begin{itemize}
% \item \textbf{Split/Split (green):} highest and smoothest composite growth.
% \item \textbf{Inert/Inert (red):} lowest growth, most disrupted by primes.
% \item \textbf{Blue and Orange:} fall between these extremes.
% \end{itemize}
% 
% This supports the hypothesis that modular prime structure governs composite background growth.
% 
% >>>>>> end of annexed block <<<<<<
\section*{Machine Learning the Composite Signals}

To explore how early arithmetic patterns might generalize, Claude \href{https://colab.research.google.com/drive/1kNbqnV0W_gIFTbSCFUpgpWQTHq2MpWZ2?usp=sharing}{trained} a \textbf{Boosted Decision Tree (BDT)} classifier on integers up to $n = 300$ using the following features:
\begin{itemize}
  \item $\Delta'(n)$: the numerical derivative of the composite log accumulation function,
  \item $n \bmod 3$: capturing Eisenstein modular class,
  \item $n \bmod 4$: capturing Gaussian modular class.
\end{itemize}

The target label was whether $n$ is an \emph{Eisenstein-inert prime}, i.e.\ $p \equiv 2 \pmod{3}$.  
The BDT was then applied to integers in the range $n = 301$ to $1000$. The classifier achieved:
\begin{itemize}
  \item \textbf{Recall}: $1.00$ — all Eisenstein-inert primes in the test range were correctly identified,
  \item \textbf{Precision}: $0.47$ — nearly half of the predicted inert primes were correct.
\end{itemize}

This shows that gradient and modular features from early integers retain predictive power in identifying structurally defined primes in unseen regions. It validates $\Delta'(n)$ and modular residue as informative, generalizable descriptors — opening the door to arithmetic learning models trained on finite regimes and extrapolated over the number line.

\begin{tcolorbox}[
  title=Why Irregular Prime Prediction is Hard,
  colback=blue!5!white,
  colframe=blue!75!black,
  boxrule=0.8pt,
  arc=6pt,
  left=6pt,right=6pt,top=6pt,bottom=6pt
]
A prime $p$ is \textbf{irregular} if it divides the numerator of at least one Bernoulli number $B_{2k}$ for $0<2k<p-1$:
\[
p \mid \mathrm{num}(B_{2k}) \quad \text{for some $2k<p-1$.}
\]
The standard method — e.g.\ the Akiyama–Tanigawa algorithm — involves computing many Bernoulli numbers modulo $p$, with complexity $O(p^{2})$ per prime. Testing all primes up to $n$ then costs roughly $O\!\bigl(n^{3}/\log n\bigr)$, making this infeasible for large $n$.
\end{tcolorbox}

\noindent\textbf{Celiac:} What are you actually feeding into the classifier?  

\noindent\textbf{Claude:} Just the essentials. Each number $n$ gets: $\Delta'(n)$, $n\bmod3$, and $n\bmod4$. These modular classes encode splitting in the Eisenstein and Gaussian integers.  

\noindent\textbf{Celiac:} And it generalizes?  

\noindent\textbf{Claude:} Surprisingly well. On Eisenstein-inert primes: recall = 1.0, precision = 0.47.  

\noindent\textbf{Jovannah:} Hmm. I didn’t expect Celiac to be impressed. But then again, you are famously supportive. Must be the politeness drilled into you.  

\noindent\textbf{Celiac:} It’s structure, Jovannah. Claude didn’t predict all primes — only the ones with clean modular roots.  

To place Claude’s shortcuts in context, consider three levels of arithmetic complexity:

\begin{tcolorbox}[
  title=The Arithmetic Complexity Staircase,
  colback=blue!2!white,
  colframe=blue!75!black,
  boxrule=0.7pt,
  arc=4pt,
  left=6pt,right=6pt,top=4pt,bottom=4pt
]
\textbf{Level 1:} \texttt{Residues} — Computing $n\bmod m$ is $O(1)$ per number. Full scan: $O(n)$.  

\vspace{0.3em}
\textbf{Level 2:} \texttt{Euler Totient} $\phi(n)$ — Sieve-based computation over $n$: $O(n\log\log n)$.  

\vspace{0.3em}
\textbf{Level 3:} \texttt{Irregular Primes via Bernoulli Numerators} — Cost: $\boxed{O(n^{3}/\log n)}$. Slowest and most accurate.  
\href{https://colab.research.google.com/drive/1K1Z1-uYhPIOCX89DaqeicSSG-k7JYkNH?usp=sharing}{[Bernoulli Reference]}
\end{tcolorbox}

% >>>>>> ANNEX C (cut-sheet v2): model benchmarking: out-of-sample precision tables — block commented out below, pointer left in text <<<<<<
\textit{(The horse-race between models --- features, folds, precision tables --- runs in Annexe C and the companion notebooks; Claude kept only the moral: fewer dials, better bets.)}

% \subsection*{Table 1: Out-of-Sample Precision of Various Models}
% 
% \begin{center}
% \begin{tabular}{@{}lcccc@{}}
% \toprule
% \textbf{Model} & \textbf{Recall (\%)} & \textbf{Precision} & \textbf{F1 Score} & \textbf{ROC AUC} \\\\ 
% \midrule
% True Algorithm (Bernoulli) & 100 & 100 & 1.00 & 1.00 \\\\
% Random Guessing & 44 & 44 & 0.44 & 0.50 \\\\
% ML Model (Full Features) & 8 & 39 & 0.14 & 0.50 \\\\
% ML Model (5 Pruned Factors) & 3 & 29 & 0.05 & 0.47 \\\\
% \bottomrule
% \end{tabular}
% \end{center}
% 
% \noindent\textbf{Jovannah:} So your classifier was outperformed fivefold by random guessing.  
% 
% \noindent\textbf{Claude:} In precision, yes. But it’s ranking, not thresholding, that helps. The top-ranked candidates have meaning.  
% 
% \noindent\textbf{Jovannah:} You’re building risk models without modeling the underlying obligor — i.e.\ Bernoulli divisibility.  
% 
% \noindent\textbf{Claude:} I knew it was heuristic. My aim was to reduce the cost — from $O(n^{3})$ to something tractable.  
% 
% \noindent\textbf{Celiac:} And you did. You found weak signal in arithmetic noise. That’s worth something.  
% 
% \noindent\textbf{Jovannah:} Perhaps. But don’t mistake weak signal for understanding.  
% 
% \noindent\textbf{Claude:} Fair. But even in markets, heuristic signal beats knowing nothing. And residues are fast.  
% 
% \noindent\textbf{Celiac:} Let’s say your model’s value isn’t precision — it’s \textit{directionality}.  
% 
% \noindent\textbf{Claude:} I’ll take that. Call it a Bayesian prior for irregularity.  
% 
% \noindent\textbf{Jovannah:} But only as long as it’s footnoted: \textit{All statistical signal here is post hoc and contingent. The Bernoulli definition remains primary.}  
% 
% 
% \section*{Machine Learning the Composite Signals}
% 
% To explore how early arithmetic patterns might generalize, Claude \href{https://colab.research.google.com/drive/1kNbqnV0W_gIFTbSCFUpgpWQTHq2MpWZ2?usp=sharing}{trained} a \textbf{Boosted Decision Tree (BDT)} classifier on integers up to $n = 300$ using the following features:
% \begin{itemize}
% \item $\Delta'(n)$: the numerical derivative of the composite log accumulation function,
% \item $n \bmod 3$: capturing Eisenstein modular class,
% \item $n \bmod 4$: capturing Gaussian modular class.
% \end{itemize}
% 
% The target label was whether $n$ is an \textit{Eisenstein-inert prime}, i.e., $p \equiv 2 \mod 3$. The BDT was then applied to integers in the range $n = 301$ to $1000$. The classifier achieved:
% \begin{itemize}
% \item \textbf{Recall}: $1.00$ --- all Eisenstein-inert primes in the test range were correctly identified,
% \item \textbf{Precision}: $0.47$ --- nearly half of the predicted inert primes were correct.
% \end{itemize}
% 
% This shows that gradient and modular features from early integers retain predictive power in identifying structurally defined primes in unseen regions. It validates $\Delta'(n)$ and modular residue as informative, generalizable descriptors --- opening the door to arithmetic learning models trained on finite regimes and extrapolated over the number line.
% 
% The study of $\Delta(n) = \log n! - \vartheta(n)$ offers a fresh signal-theoretic perspective on the structure of composite integers, exposing patterns hidden beneath the prime scaffolding. From Euler's and Stirling's expansions to spectral entropy and machine learning, this approach unifies classic analysis with modern interpretive tools.
% 
% The modular prime structure --- particularly under $(\bmod 4, \bmod 3)$ congruences --- leaves detectable signatures in the composite background. The smoothed growth rates $\Delta'_{i,j}(n)$ reflect how arithmetic structure shapes the "roughness" of integer accumulation. Furthermore, these signals prove learnable: classifiers trained on local patterns can generalize modular behavior into broader regions.
% 
% Together, this suggests a new framing of integer complexity: not merely in the distribution of primes, but in the information-theoretic geometry of what fills the gaps --- the entropic landscape of the composites.
% 
% %%%%%%%%%%%%%%%%
% 
% “What are you actually feeding into the classifier?” Celiac asked, peering at the scroll of Python code.
% 
% Claude grinned. “Only the essentials,” he replied. “Each number $n$ gets three ingredients: the local gradient $\Delta'(n)$ of the composite log excess, and its modular residues $n \bmod 3$ and $n \bmod 4$. These modular classes encode how $n$ interacts with the Eisenstein and Gaussian integers — essentially, whether it’s inert or split in those rings.”
% 
% He continued, “We train a Boosted Decision Tree (BDT) on just the integers up to $n = 300$. The goal? Learn to recognize Eisenstein-inert primes — those $p \equiv 2 \mod 3$ — purely from these features. Then, we test the classifier on numbers from $301$ to $1000$.”
% 
% Celiac raised an eyebrow. “And it works?”
% 
% “It works surprisingly well,” Claude nodded. “The model achieves perfect \textbf{recall} — it flags all the Eisenstein-inert primes. And nearly half the time, it’s right — \textbf{precision} around $0.47$. Potentially impressive considering the simplicity and short training window.”
% 
% He opened a notebook on his screen and gestured at the table:
% 
% \begin{center}
% \begin{tabular}{|c|c|c|c|}
% \hline
% Index & Prime & Mean Predicted Probability & Std Dev \\
% \hline
% 820 & 17891 & 0.000893 & 0.000278 \\
% 1003 & 19727 & 0.000893 & 0.000278 \\
% 904 & 18691 & 0.000893 & 0.000278 \\
% 592 & 15629 & 0.000889 & 0.000272 \\
% ... & ... & ... & ... \\
% \hline
% \end{tabular}
% \end{center}
% 
% “These are the top candidates for irregular primes in the range $10,001$ to $20,000$, identified using an ensemble of BDTs with dropout-style uncertainty estimation,” Claude explained. “The mean probabilities are low because the task is harder, but I believe the ranking is stable.”
% 
% \vspace{0.5em}
% Celiac studied the predictions. “So even the background accumulation signal, your $\Delta'(n)$, knows something about prime inertia?”
% 
% “Yes,” Claude said, “and this is the deeper point. We’ve been treating $\Delta(n) = \log n! - \vartheta(n)$ as a kind of signal: its derivative captures how the composites ‘fill in’ the logarithmic mass. But what we now see is that this signal is not only measurable: it’s learnable.”
% 
% He paused. “The modular prime structure: especially in $(\bmod 4, \bmod 3)$ congruence classes, leaves subtle but detectable imprints in the background composite growth. The smoothed curves $\Delta'_{i,j}(n)$ for each class capture these differences in slope. And the BDT is picking up on them.”
% 
% \vspace{1em}
% \noindent
% \subsection*{Predicting Irregular Primes via Boosted Decision Trees}
% 
% Claude was of the view that he was glimpsing a complementary picture: an information-theoretic geometry not of the spikes (the primes), but of what lies between them: the entropic landscape of the composites.
% 
% \begin{tcolorbox}[
%   title=Irregular Primes via Bernoulli Numbers,
%   colback=blue!5!white,
%   colframe=blue!75!black,
%   boxrule=0.8pt,
%   arc=6pt,
%   left=6pt,right=6pt,top=6pt,bottom=6pt
% ]
% A prime \(p\) is \emph{irregular} if it divides the numerator of at least one Bernoulli number \(B_{2k}\) with \(0<2k<p-1\):  
% \[
% p \mid \mathrm{num}(B_{2k}) \quad \text{for some } 0<2k<p-1.
% \]
% This is the only mathematically correct criterion. All other methods, including the model below, are heuristics that mimic credit‑risk scoring without knowing the true obligor balance sheet.
% \end{tcolorbox}
% 
% To investigate whether the signature of {\em irregularity} might be learnable from some of these humble classical arithmetic features, Claude trained a \href{https://colab.research.google.com/drive/1kNbqnV0W_gIFTbSCFUpgpWQTHq2MpWZ2?usp=sharing}{Boosted Decision Tree (BDT)} classifier using known irregular primes up to 10,000. Each prime was represented by the following features:
% \begin{itemize}
%     \item $\log p$: logarithmic scale of the prime,
%     \item $\psi(p) + \gamma - \log p$: deviation from harmonic prediction,
%     \item $\frac{d}{dp}\vartheta(p)$: local Chebyshev slope,
%     \item Modulo residues: $p \bmod 3$, $p \bmod 4$, $p \bmod 5$, $p \bmod 6$, $p \bmod 7$, $p \bmod 8$.
% \end{itemize}
% 
% The trained model was then applied to primes in the range $5,\!001$ to $10,\!000$. The classifier was able to generate a ranked list of the most likely candidates for irregularity, suggestive that weak signals of arithmetic anomaly might persist at higher primes.
% 
% \begin{center}
% \begin{tabular}{|c|c|c|}
% \hline
% \textbf{Rank} & \textbf{Prime} & \textbf{Predicted Probability} \\
% \hline
% 1 & 18503 & 0.000225 \\
% 2 & 19727 & 0.000193 \\
% 3 & 19937 & 0.000193 \\
% 4 & 19319 & 0.000190 \\
% 5 & 10957 & 0.000187 \\
% 6 & 12721 & 0.000187 \\
% 7 & 13781 & 0.000182 \\
% 8 & 16529 & 0.000175 \\
% \hline
% \end{tabular}
% \end{center}
% 
% Even with \href{https://colab.research.google.com/drive/1qe-CSkkfjqBhuJaffiqEXhSMNDF1OdCY?usp=sharing}{enhanced} features and deeper trees, the predicted probabilities remained below $0.03\%$ suggesting a ranking of candidates. 
% 
% 
% 
% \vspace{1em}
% 
% \noindent
% \textbf{Claude}  
% We built a five‑factor model on primes up to 7000. It had “alpha”: F1 of 0.41 in‑sample, 44\% recall. We pruned “leaky” features like cumulative irregular density:just like a risk desk removing lagged default rates.  
% 
% \noindent
% \textbf{Table 1. Out-of-Sample Performance on Primes 7001--10000}
% 
% \begin{center}
% \begin{tabular}{@{}lcccc@{}}
% \toprule
% \textbf{Model} & \textbf{Recall (\%)} & \textbf{Precision} & \textbf{F1 Score} & \textbf{ROC AUC} \\
% \midrule
% Akiyama--Tanigawa Algo. & 100 & 100 & 1.00 & 1.0 \\
% Random Guessing (base) & 44 & 44 & 0.44 & 0.5 \\
% ML Model (18 features) & 8 & 39 & 0.14 & 0.5 \\
% ML Model (5 pruned factors) & 3 & 29 & 0.05 & 0.5 \\
% \bottomrule
% \end{tabular}
% \end{center}
% \noindent
% \textbf{Jovannah:}  
% And out‑of‑sample?  
% 
% \noindent
% \textbf{Claude:}  
% Recall collapsed to 2.8\%. Precision 28.6\%. AUC 0.47. Worse than random.  
% 
% \noindent
% \textbf{Jovannah:}  
% So you built a subprime CDO model for primes.  
% You \href{https://colab.research.google.com/drive/1K1Z1-uYhPIOCX89DaqeicSSG-k7JYkNH?usp=sharing}{removed} the only feature with signal: “historical default rate” a.k.a.\ cum\_irreg\_density, and now you’re shocked there’s no alpha left?  
% 
% \noindent
% \textbf{Claude:}  
% But we still had five structural “risk factors”:
% \begin{itemize}
% \item \(\mathrm{radical}(p-1)\) (14.8\%) ,  normalized \(\phi(p-1)\) (13.9\%)  
% \item gap ratio (9.4\%) , harmonic deviation (6.2\%) , \(\sqrt{p}\) (6.6\%)  
% \end{itemize}
% They looked like a Fama–French for primes.  
% 
% \noindent
% \textbf{Jovannah:}  
% Looked, yes. But your factors had no link to the underlying credit event—Bernoulli number divisibility.  
% You were modelling loyalty card expenditure to predict mortgage defaults. In 2006 it would have backtested just fine. In 2008 you'd gone bankrupt like the best of them.  
% 
% \noindent
% \textbf{Claude:}  
% Agreed, that even our 18‑factor model with cum\_irreg\_density was effectively just memorising base rates?  
% 
% \noindent
% \textbf{Jovannah:} Exactly. Like a VAR model trained on low‑vol years.  
% In-sample it looks clever. Out-of-sample it dies.  
% Irregular primes aren’t “noisy obligors” you can score, rather they’re deterministically defined by Bernoulli numerators.  
% 
% \vspace{0.5em}
% 
% \noindent
% \textbf{Claude:} I knew this, of course. It was never meant to rival the Bernoulli test.  
% Just a heuristic: to cut down\footnote{When you begin to think like a quant, you ask as Claude did: “How much structure can I extract without paying cubic cost?”   The answer may well be: residues.   For any prime $p$, computing $p \bmod m$ is constant time.  Do this for $n$ primes, you get an $O(n)$ complexity pipeline.  So he threw in \texttt{mod3, mod4, mod6} as predictors; not because they were brilliant, but because they were cheap. And sometimes, cheap structure is all you need to beat randomness.} on the $O(n^3)$ computational slog\footnote{
% Claude's model, though naive in the face of arithmetic determinism, serves a pragmatic role. The Bernoulli test (via Akiyama–Tanigawa or von Staudt–Clausen) scales poorly: testing a prime $p$ for irregularity involves examining roughly $\lfloor (p - 3)/2 \rfloor$ Bernoulli numerators. Each numerator is nontrivial to compute and costly modulo large $p$. Hence the appeal of a fast, learnable proxy. Yet, as Jovannah rightly points out, this approach trades certainty for signal; and when the signal is weak, even elegant models collapse under precision scrutiny.} of checking divisibility across all $B_{2k}$.
% 
% \noindent
% \textbf{Jovannah:} \\
% The same journey as every junior quant:
% \begin{enumerate}
%   \item In‑sample optimism (“we have alpha!”).  
%   \item Out‑of‑sample collapse (“we have noise”).  
%   \item And finally: Model the mechanism, not the shadow.
% \end{enumerate}
% \noindent
% 
% Your Table~1 says it all: even random guessing beats your pruned model fivefold, and the true algorithm is perfect. 
% 
% \noindent
% \textbf{Claude:}  
% I just wanted to avoid computing \( B_{2k} \) up to size 5000 — too expensive.  
% With the totient function for example, you can sieve all values up to 10,000 in less than a second:
% 
% \begin{lstlisting}[language=Python, caption=Euler's Totient Function up to a Limit, label=lst:totient, basicstyle=\ttfamily\small]
% def compute_totients(limit):
%     phi = list(range(limit + 1))  # phi[n] = n initially
% 
%     for i in range(2, limit + 1):
%         if phi[i] == i:  # i is prime
%             for j in range(i, limit + 1, i):
%                 phi[j] -= phi[j] // i
% 
%     return phi
% \end{lstlisting}
% 
% \noindent
% \textbf{Jovannah:}  
% Exactly. That's because \( \phi(n) \) is multiplicative and lends itself to a sieve. But you’re now comparing a shallow slope to a vertical wall.  
% Sieve of Eratosthenes gives you \( O(n \log \log n) \),  
% But computing irregularity is a \( O(n^3 / \log n) \) business.
% 
% \noindent
% \textbf{Claude:}  
% So until we can model Bernoulli numerators mod \(p\), no amount of “factor engineering” will work?
% 
% \noindent
% \textbf{Jovannah:}  
% Precisely. You’re trying to forecast a deterministic number‑theoretic property with cosmetic risk factors.  
% It’s an instructive case study in overfitting and false confidence—but not a path to irregular prime discovery.  
% 
% \noindent
% \begin{tcolorbox}[
%   title=Why Bernoulli Irregularity Testing is Expensive,
%   colback=blue!5!white,
%   colframe=blue!75!black,
%   boxrule=0.8pt,
%   arc=6pt,
%   left=6pt,right=6pt,top=6pt,bottom=6pt
% ]
% To test whether a prime \( p \) is irregular, one must compute all Bernoulli numbers \( B_{2k} \) for \( 2 \leq 2k \leq p - 3 \), and check if \( p \mid \text{numerator}(B_{2k}) \).  
% Using optimized modular algorithms (e.g. Akiyama–Tanigawa), the cost per prime is \( O(p^2) \).  
% Hence testing all primes up to \( n \) costs approximately:
% 
% \[
% \boxed{O(n^3 / \log n)}
% \]
% 
% That is orders of magnitude beyond what a quick ML heuristic offers.
% \end{tcolorbox}
% 
% 
%%%%%%%%%%%%%%%%%%%%
% 
% \section*{Motivating Composite Logarithmic Roughness}
% 
% \subsection*{1. Derivation of the Logarithmic Derivative of $n!$}
% 
% Euler provided a foundational insight by recognizing that the logarithm of the factorial, though a discrete object, can be interpolated and differentiated using analytic tools. Begin with the identity:
% \begin{equation*}
% \log n! = \sum_{k=1}^n \log k,
% \end{equation*}
% which expresses $\log n!$ as a sum over the natural logarithms of all integers up to $n$.
% 
% To analyze this with continuous tools, we approximate the discrete sum by an integral. This motivates a smooth interpolation of the factorial via the gamma function, where $n! = \Gamma(n+1)$ and thus:
% \begin{equation*}
% \log n! = \log \Gamma(n+1).
% \end{equation*}
% 
% Differentiating both sides with respect to $n$ gives:
% \begin{equation*}
% \frac{d}{dn} \log n! = \frac{d}{dn} \log \Gamma(n+1) = \psi(n+1),
% \end{equation*}
% where $\psi(x)$ is the digamma function — the logarithmic derivative of the gamma function:
% \begin{equation*}
% \psi(x) = \frac{d}{dx} \log \Gamma(x).
% \end{equation*}
% 
% Euler himself studied this function in relation to harmonic numbers. For integer values, we have the exact identity:
% \begin{equation*}
% \psi(n+1) = H_n - \gamma,
% \end{equation*}
% where $H_n = \sum_{k=1}^n \frac{1}{k}$ is the $n$th harmonic number, and $\gamma \approx 0.5772$ is the Euler–Mascheroni constant.
% 
% \medskip
% 
% \noindent
% So, in summary:
% \begin{equation*}
% \frac{d}{dn} \log n! = \psi(n+1) = H_n - \gamma.
% \end{equation*}
% 
% This elegant result connects factorial growth — often considered “multiplicative” — to the additive structure of harmonic sums. It highlights how the accumulation of $\log k$ behaves like the accumulation of $1/k$, and serves as a cornerstone in both analytic number theory and information theory.
% 
% \subsection*{2. The Chebyshev Function $\vartheta(n)$ and Prime Contribution}
% 
% We define the \textbf{primorial function} as:
% \begin{equation}
% p_{\#} = \prod_{\substack{p \leq n \ p ,\text{prime}}} p,
% \end{equation}
% and its logarithm is the Chebyshev function:
% \begin{equation}
% \vartheta(n) = \log(p(_{\#}) = \sum_{\substack{p \leq n \ p,\text{prime}}} \log p.
% \end{equation}
% 
% Unlike $\log n!$, which accumulates the logarithms of \emph{all} integers up to $n$, $\vartheta(n)$ sums only over the primes. As a result, $\vartheta(n)$ increases in discrete jumps at each prime, while $\log n!$ is a smooth, cumulative function.
% 
% The difference between them defines the composite logarithmic excess:
% \begin{equation}
% \Delta(n) = \log n! - \vartheta(n)
% \end{equation}
% This quantity measures the total log-volume added by \emph{composite} integers up to $n$.
% 
% \begin{figure}[h!]
% \centering
% \includegraphics[width=0.4\textwidth]{Illustrations/PuboidFactorials.png}
% \caption{Log-volume interpretation: $\log n!$ is decomposed into the prime contribution $\vartheta(n)$ and the composite remainder $\Delta(n)$.}
% \end{figure}
% 
% \subsection*{3. Composite Roughness as an Entropic Signal}
% 
% The function $\Delta(n)$ exhibits an intriguing behavior: it grows smoothly between primes, with dips aligned at prime locations where $\vartheta(n)$ jumps. When plotted, its derivative appears jagged due to these prime-induced interruptions. By smoothing the derivative (e.g., via moving averages), we obtain a signal that reflects the overall rhythm by which composites fill the gaps between primes.
% 
% This smoothed derivative can be interpreted as a kind of \textbf{composite roughness measure} — analogous to entropy in statistical mechanics or information theory. Like the Hurst exponent, which characterizes long-term memory in time series, or Shannon entropy, which quantifies unpredictability, this measure reflects the structure and distribution of primes versus composites.
% 
% To make the Hurst analogy more concrete, we note that the signal $\Delta'(n)$ exhibits fluctuating but persistent trends in the local rate of logarithmic accumulation by composites. If we treat $\Delta'(n)$ as a time series, we can apply rescaled range analysis or compute the variance of cumulative deviations to assess whether the signal has long-range dependence. A Hurst exponent $H > 0.5$ would indicate persistent behavior — that is, intervals rich in composite contributions are likely to be followed by similarly rich intervals, suggesting structured filling of gaps between primes.
% 
% Moreover, in signal-to-noise terms, $\vartheta(n)$ acts like impulsive noise — introducing high-frequency jumps — while $\log n!$ is a low-frequency baseline. The difference $\Delta(n)$ then encodes the signal formed by the superposition of compositeness over prime-induced disruptions. The roughness of this signal captures a kind of arithmetic entropy — a measure of the unpredictability or smoothness of integer structure beneath the prime scaffolding.
% 
% To make the connection to Shannon entropy more precise, we consider the distribution of contributions to $\Delta(n)$ from each composite $k \leq n$ as a probability mass function:
% \begin{equation}
% P(k) = \frac{\log k}{\sum_{j \in \text{composites} \leq n} \log j}
% \end{equation}
% 
% Then the Shannon entropy of this distribution is:
% \begin{equation}
% S(n) = - \sum_{k \in \text{composites} \leq n} P(k) \log P(k)
% \end{equation}
% 
% This entropy $S(n)$ measures the information-theoretic complexity of the composite structure beneath $n$, reflecting how evenly or unevenly the logarithmic mass is distributed among composites. A more uniform distribution yields higher entropy, suggesting richer structural variety, while a skewed distribution reflects more concentrated contributions.
% 
% The growth of $S(n)$ suggests a gradual enrichment in the diversity of logarithmic contributions from composites as $n$ increases. At $n=1000$, the entropy reaches approximately $6.71$, a strong indicator of the nuanced structure in the distribution of composite integers.
% 
% To complement this, we compute the \textbf{spectral flatness} of the smoothed derivative $\Delta'(n)$. This metric compares the geometric mean of the power spectrum to its arithmetic mean. A value near 1 indicates white-noise-like behavior; a value near 0 indicates tonal or highly structured content. We find:
% \begin{equation}
% \text{Spectral Flatness} \approx 0.309
% \end{equation}
% 
% This moderate flatness supports the interpretation that composite roughness is neither pure noise nor fully periodic but exhibits structured, quasi-random behavior.
% 
% \paragraph{Future signal-to-noise perspectives.} Additional analysis could include:
% \begin{itemize}
% \item \textbf{Autocorrelation decay:} quantify periodicity and memory.
% \item \textbf{Signal-to-noise ratio (SNR):} via decomposition into trend and fluctuation.
% \item \textbf{Block entropy or entropy rate:} proxy for Kolmogorov complexity.
% \item \textbf{Cross-spectral coherence:} compare with prime indicator sequences.
% \end{itemize}
% 
% 
%%%%%%%%%%%%%%%%%%
% 
% >>>>>> end of annexed block <<<<<<
% >>>>>> ANNEX E (cut-sheet v2): benchmarking qubit sufficiency by Shor's algorithm — block commented out below, pointer left in text <<<<<<
\textit{(How many qubits it takes to break the world's locks --- Shor's algorithm benchmarked step by step --- is set out in Annexe E.)}

% \section*{Benchmarking Qubit Sufficiency by Shor's Algorithm}
% 
% Quantum computing promises exponential speedups for certain problems, and among them, \textbf{Shor's algorithm} for integer factorization remains the most iconic. More than a cryptographic threat, Shor's algorithm is a \textit{yardstick} for assessing when quantum hardware becomes truly useful: it tells us how many qubits—both physical and logical—are sufficient to tackle problems beyond classical reach.
% 
% \vspace{1em}
% \noindent
% A quantum computer with \( n \) qubits can theoretically store a superposition over \( 2^n \) amplitudes. For example:
% \[
% 2^{100} \approx 1.27 \times 10^{30}
% \]
% This represents a Hilbert space of unimaginable size—far beyond the capabilities of any classical machine. Yet Shor's algorithm is not merely concerned with raw state space, but with the \textit{structured use of qubits} in register-based computation. To factor an \( n \)-bit number \( N \), Shor’s algorithm typically requires:
% \begin{itemize}
%   \item A \textit{control register} of size \( 2n + \log n \) qubits, for quantum phase estimation
%   \item A \textit{work register} of \( n \) qubits, for modular arithmetic and exponentiation
%   \item Approximately \( n \) \textit{ancilla qubits}, for arithmetic carry and temporary storage
% \end{itemize}
% A common rule of thumb is then \(
% \text{Total logical qubits required} \approx 4n.
% \)
% 
% \begin{table}[h!]
% \centering
% \begin{tabular}{|c|c|c|}
% \hline
% RSA Key Size & \( n \) (bits) & Estimated Logical Qubits Needed \\
% \hline
% 512-bit RSA  & 512   & \(\approx 2,\!048\) \\
% 1024-bit RSA & 1024  & \(\approx 4,\!096\) \\
% \hline
% \end{tabular}
% \caption{Approximate logical qubit requirements for Shor's algorithm}
% \end{table}
% 
% \subsection*{Physical Qubit Requirements (with QEC)}
% 
% Due to decoherence and gate noise, logical qubits must be protected by \textbf{quantum error correction} (QEC), often using surface codes. A typical overhead: \(
% \text{1 logical qubit} \Rightarrow \text{1000 physical qubits}
% \)
% This implies the following physical qubit budgets: RSA-512: ~2 million physical qubits, RSA-1024: ~4 million, highlighting the scale required to break classical encryption — and serve as a benchmark for hardware design and national security. While Shor's algorithm is traditionally formulated using \textit{qubits} (2-level systems), its structure is generalizable. Using \textbf{qutrits} (3-level systems) or higher-dimensional \textit{qudits} may:
% \begin{itemize}
%   \item Reduce circuit depth by encoding more information per element
%   \item Allow more compact arithmetic and parallelism in modular operations
%   \item Offer new forms of error resilience via richer algebraic structure
% \end{itemize}
% If quantum error correction can be extended to qutrits with competitive overheads, the effective resource cost of Shor’s algorithm—and many others—could be significantly reduced. In this sense, Shor's algorithm remains not only a benchmark, but a gateway to exploring how \textit{dimensionality} impacts quantum logic.
% 
% \section*{Finding the Period for \boldmath$a = 13$ Modulo $N = 1260$}
% We'll illustrate the key steps of Shor’s algorithm on the composite that Demi would have us know is the product of two \textit{puboids} \( N = 1260 \), with a particular focus on the modular periodicity that underlies its power.:
% \[
% N = 1260 = 2^2 \cdot 3^2 \cdot 5 \cdot 7
% \]
% The algorithm requires us to choose an integer \( a \) such that:
% \[
% 1 < a < N \quad \text{and} \quad \gcd(a, N) = 1
% \]
% This ensures \( a \) belongs to the multiplicative group \( \mathbb{Z}_N^* \), allowing modular exponentiation to define a meaningful cyclic structure. We choose:
% \[
% a = 13
% \]
% This value is not one of the prime factors of 1260, and so satisfies:
% \[
% \gcd(13, 1260) = 1
% \]
% so it is a valid candidate. The next step is to determine the \textit{period}, of \( a \) modulo \( N \). That is, we seek the smallest positive integer \( r \) such that:
% \[
% 13^r \equiv 1 \mod 1260
% \]
% This is the step that Shor's algorithm performs efficiently using quantum phase estimation. Here, we simulate it classically to illustrate the behavior. Below is the table of successive powers of 13 modulo 1260:
% 
% \begin{center}
% \begin{tabular}{|c|c|}
% \hline
% Exponent \( k \) & \( 13^k \mod 1260 \) \\
% \hline
% 1  & 13   \\
% 2  & 169  \\
% 3  & 877  \\
% 4  & 91   \\
% 5  & 943  \\
% 6  & 889  \\
% 7  & 217  \\
% 8  & 301  \\
% 9  & 1213 \\
% 10 & 529  \\
% 11 & 577  \\
% 12 & 1    \\
% \hline
% \end{tabular}
% \end{center}
% We observe that:
% \[
% 13^{12} \equiv 1 \mod 1260
% \]
% and no smaller exponent satisfies this, so the \textit{period is}:
% \(
% \boxed{r = 12}.
% \)
% This periodicity underpins the core quantum speedup in Shor’s algorithm. The next step is to use this period to extract a nontrivial factor of \( N \), which is implemented in \href{https://colab.research.google.com/drive/1hn_NBpzrlv1TNKeRQAPDLabW04Jya0tv?usp=sharing}{Python Implementation} as follows
% 
% \begin{lstlisting}[language=Python]
% import pandas as pd
% import math
% from math import gcd
% 
% def find_period_verbose(a, N):
%     if gcd(a, N) != 1:
%         return [], None, f"gcd({a}, {N}) = {gcd(a, N)} → already a nontrivial factor!"
%     results = []
%     current = 1
%     for k in range(1, 1001):
%         current = (current * a) % N
%         results.append((k, current))
%         if current == 1:
%             return results, k, None
%     return results, None, "Period not found."
% 
% def shor_post_processing(a, r, N):
%     a_r2 = pow(a, r // 2, N)
%     plus = (a_r2 + 1) % N
%     minus = (a_r2 - 1) % N
%     return a_r2, plus, minus, gcd(plus, N), gcd(minus, N)
% 
% def lcm(x, y):
%     return x * y // gcd(x, y)
% 
% def reconstruct_from_factors(f1, f2):
%     return lcm(f1, f2)
% \end{lstlisting}
% \newpage
% 
% \section*{Run for \boldmath\( N = 1260 \), \( a = 13 \)}
% 
% \begin{lstlisting}[language=Python]
% a = 13
% N = 1260
% 
% results, period, note = find_period_verbose(a, N)
% df = pd.DataFrame(results, columns=["Exponent k", f"{a}^k mod {N}"])
% display(df)
% 
% if note:
%     print(note)
% else:
%     a_r2, plus, minus, gcd_plus, gcd_minus = shor_post_processing(a, period, N)
%     reconstructed_N = reconstruct_from_factors(gcd_plus, gcd_minus)
% 
%     print(f"a^(r/2) mod N = {a_r2}")
%     print(f"a^(r/2) + 1 = {plus}, gcd = {gcd_plus}")
%     print(f"a^(r/2) - 1 = {minus}, gcd = {gcd_minus}")
%     print(f"Reconstructed N from LCM = {reconstructed_N}")
% \end{lstlisting}
% 
% \subsection*{Example Output}
% 
% \begin{itemize}
%   \item \( r = 12 \)
%   \item \( a^{6} \equiv 1009 \mod 1260 \)
%   \item \( \gcd(1009 - 1, 1260) = 252 \), \( \gcd(1009 + 1, 1260) = 10 \)
%   \item \( \mathrm{lcm}(252, 10) = 1260 = N \)
% \end{itemize}
% 
% 
% Even if the GCDs found are not prime (10 and 252), we can still reconstruct the original composite number \( N \) via their LCM. Shor’s algorithm thus provides a quantum-assisted path to classical factor recovery.
% 
% 
% 
% 
% 
%%%%%%%%%%%%%%%%%%%%%%%%%%%%%%%%%%
% (The Green Ethical Tranche chapter moved to 8-14greenTranche.tex — cut-sheet v2)
% >>>>>> end of annexed block <<<<<<
